% Pertemuan 11 Praktikum Pemrograman (IF25-11002). Dihasilkan oleh tools/build.py dari tools/konten/p11.py.
% Kompilasi: xelatex P11.tex (dua kali). Catatan: xelatex -jobname=P11-catatan "\def\catatan{1}% Pertemuan 11 Praktikum Pemrograman (IF25-11002). Dihasilkan oleh tools/build.py dari tools/konten/p11.py.
% Kompilasi: xelatex P11.tex (dua kali). Catatan: xelatex -jobname=P11-catatan "\def\catatan{1}% Pertemuan 11 Praktikum Pemrograman (IF25-11002). Dihasilkan oleh tools/build.py dari tools/konten/p11.py.
% Kompilasi: xelatex P11.tex (dua kali). Catatan: xelatex -jobname=P11-catatan "\def\catatan{1}% Pertemuan 11 Praktikum Pemrograman (IF25-11002). Dihasilkan oleh tools/build.py dari tools/konten/p11.py.
% Kompilasi: xelatex P11.tex (dua kali). Catatan: xelatex -jobname=P11-catatan "\def\catatan{1}\input{P11.tex}"
\documentclass[t]{beamer}
\usetheme{ITERA}
\input{catatan-preamble.tex}
\renewcommand\ITERAmeeting{Pertemuan 11}
\begin{document}
\SlideJudul{IF25-11002 PRAKTIKUM PEMROGRAMAN}{Pertemuan 11\\Searching}{Mencari data di array dengan linear search dan binary search, dan membandingkan banyak langkahnya}{Tim Pengajar}{Tahun 2026. Sejajar dengan Pertemuan 11 Algoritma Pemrograman: Searching.}
\note{Buka dengan menanyakan satu hal dari pertemuan lalu yang masih membingungkan.}

\begin{frame}{Dua mata kuliah, satu minggu}
\framesubtitle{Sebelum mulai}
Topik minggu ini sama di dua mata kuliah, dengan peran yang berbeda.
\vspace{10pt}

\begin{columns}[T]
\begin{column}{0.47\textwidth}
\begin{Blok}{Algoritma: Searching}
Algoritma pencarian beruntun dan biner, prasyarat data terurut, serta menghitung banyak perbandingan.
\end{Blok}
\end{column}
\begin{column}{0.47\textwidth}
\begin{BlokGelap}{Praktikum: Searching}
Menulis linear search dan binary search di C++, menelusuri low, high, dan mid, serta menghitung perbandingan.
\end{BlokGelap}
\end{column}
\end{columns}
\note{Tanyakan apa yang sudah dibahas di Algoritma minggu ini. Gunakan jawabannya sebagai jembatan ke praktikum.}
\end{frame}

\begin{frame}{Saat pulang nanti, kalian bisa}
\framesubtitle{Target hari ini}
\begin{columns}[T]
\begin{column}{0.47\textwidth}
\begin{Langkah}{1}{Menerapkan}
Menulis program yang mencari data pada array dengan linear search dan binary search, termasuk mencari semua kemunculan dan menangani data yang tidak ditemukan

{\footnotesize\color{ITERAInkMuted}Sub-CPMK 11.1, CPMK0607}
\end{Langkah}
\end{column}
\begin{column}{0.47\textwidth}
\begin{Langkah}{2}{Menjelaskan}
Menelusuri langkah binary search dengan trace table low, high, dan mid, serta menjelaskan perbedaan banyak perbandingan kedua algoritma

{\footnotesize\color{ITERAInkMuted}Sub-CPMK 11.2, CPMK0608}
\end{Langkah}
\end{column}
\end{columns}
\vspace{8pt}
\begin{Catatan}
Dua kemampuan ini dinilai sama berat di Exit Ticket, tugas, UTS, dan UAS.
\end{Catatan}
\note{Bacakan target dengan pelan. Kata kuncinya menerapkan dan menjelaskan.}
\end{frame}

\begin{frame}{Ingat minggu lalu}
\framesubtitle{Review, 5 menit}
\begin{columns}[T]
\begin{column}{0.31\textwidth}
\begin{BlokGelap}{Pertanyaan 1}
Bagaimana menghitung rata-rata setiap kolom array 2D?
\end{BlokGelap}
\end{column}
\begin{column}{0.31\textwidth}
\begin{BlokGelap}{Pertanyaan 2}
Apa elemen diagonal utama matriks?
\end{BlokGelap}
\end{column}
\begin{column}{0.31\textwidth}
\begin{BlokGelap}{Pertanyaan 3}
Apa hasil \texttt{s.find("tik")} untuk \texttt{"Praktikum"}?
\end{BlokGelap}
\end{column}
\end{columns}
\vspace{8pt}
Jawab di kertas dulu, lalu cocokkan dengan teman sebelah.\note{Pilih dua mahasiswa untuk menjawab. Koreksi singkat, jangan mengajar ulang.}
\end{frame}

\Bagian{1}{Masalah pembuka}{Mencari nama di daftar wisuda}
\note{Masalah pembuka 10 menit.}

\begin{frame}[fragile]{Mencari nama di daftar wisuda}
\framesubtitle{Masalah minggu ini}
\begin{columns}[T]
\begin{column}{0.55\textwidth}
{Panitia wisuda memegang daftar 1.200 nama yang sudah urut abjad. Seorang orang tua bertanya, apakah anaknya ada di daftar?

\vspace{6pt}
Membaca dari nama pertama sampai ketemu bisa butuh ratusan kali melihat. Padahal saat mencari kata di kamus, kita tidak membaca dari halaman pertama.\par}
\vspace{6pt}
\begin{Catatan}
\textbf{Diskusi 2 menit.} Bagaimana kalian mencari kata di kamus? Mengapa cara itu tidak bisa dipakai pada daftar yang belum urut?
\end{Catatan}
\end{column}
\begin{column}{0.41\textwidth}
\begin{Keluaran}[]
Cari: Siti Rahma
Ditemukan di nomor 1023
Linear search : 1023 perbandingan
Binary search : 10 perbandingan
\end{Keluaran}
\end{column}
\end{columns}
\note{Tampung beberapa jawaban tanpa menilai benar salah.}
\end{frame}

\begin{frame}{Dari langkah ke instruksi}
\framesubtitle{Sambungan dengan Algoritma}
\begin{columns}[T]
\begin{column}{0.31\textwidth}
\begin{Langkah}{1}{Bandingkan satu per satu}
Cara paling sederhana: periksa setiap elemen dari depan.
\end{Langkah}
\end{column}
\begin{column}{0.31\textwidth}
\begin{Langkah}{2}{Manfaatkan urutan}
Bila data urut, buang separuh daftar di setiap langkah.
\end{Langkah}
\end{column}
\begin{column}{0.31\textwidth}
\begin{Langkah}{3}{Terjemahkan}
Linear search dengan \texttt{for}, binary search dengan \texttt{while} dan tiga indeks.
\end{Langkah}
\end{column}
\end{columns}
\vspace{10pt}
Langkah 1 dan 2 dilatih di Algoritma. Langkah 3 kita kerjakan hari ini dengan C++.\note{Hubungkan eksplisit ke Algoritma minggu ini.}
\end{frame}

\Bagian{2}{Coba dulu, baru dijelaskan}{25 menit eksperimen di GDB Online}
\note{Struggle 25 menit. Berkeliling, jangan menjelaskan materi dulu.}

\begin{frame}{Misi kalian}
\framesubtitle{Struggle, 25 menit}
\begin{columns}[T]
\begin{column}{0.31\textwidth}
\begin{Langkah}{1}{Cari}
Diberikan \texttt{int a[8] = \{4, 9, 2, 7, 5, 9, 1, 8\};}, tampilkan indeks angka 7.
\end{Langkah}
\end{column}
\begin{column}{0.31\textwidth}
\begin{Langkah}{2}{Semua}
Tampilkan semua indeks angka 9, bukan hanya yang pertama.
\end{Langkah}
\end{column}
\begin{column}{0.31\textwidth}
\begin{Langkah}{3}{Tidak ada}
Apa yang ditampilkan program kalian bila mencari angka 6?
\end{Langkah}
\end{column}
\end{columns}
\vspace{8pt}
\begin{columns}[T]
\begin{column}{0.47\textwidth}
\begin{Blok}{Boleh}
Bertanya ke teman, mengubah apa saja, sengaja membuat error.
\end{Blok}
\end{column}
\begin{column}{0.47\textwidth}
\begin{BlokAlert}{Belum dulu}
Mencari jawaban jadi di internet atau AI.
\end{BlokAlert}
\end{column}
\end{columns}
\note{Catat kesalahan yang sering muncul untuk dibahas di debrief.}
\end{frame}

\begin{frame}{Apa yang kalian temukan?}
\framesubtitle{Debrief struggle}
\begin{columns}[T]
\begin{column}{0.31\textwidth}
\begin{BlokGelap}{Pertanyaan 1}
Bagaimana program tahu sebuah data tidak ditemukan?
\end{BlokGelap}
\end{column}
\begin{column}{0.31\textwidth}
\begin{BlokGelap}{Pertanyaan 2}
Kapan \texttt{break} dipakai dan kapan tidak?
\end{BlokGelap}
\end{column}
\begin{column}{0.31\textwidth}
\begin{BlokGelap}{Pertanyaan 3}
Berapa perbandingan untuk mencari elemen terakhir?
\end{BlokGelap}
\end{column}
\end{columns}
\vspace{10pt}
Jawaban kalian akan kita uji di bagian berikutnya.\note{Tulis jawaban di papan; buktikan atau patahkan di bagian materi.}
\end{frame}

\Bagian{3}{Mencari dengan cerdas}{Linear search, binary search, dan menghitung langkah}
\note{Materi 40 menit. Tunjukkan setiap contoh langsung di GDB Online.}

\begin{frame}[fragile]{Linear search}
\framesubtitle{Periksa satu per satu}
\begin{columns}[T]
\begin{column}{0.55\textwidth}
\begin{Kode}[]{linear.cpp}
int a[8] = {4, 9, 2, 7, 5, 9, 1, 8};
int x = 7, posisi = -1;
for (int i = 0; i < 8; i++) {
    if (a[i] == x) {
        posisi = i;
        break;
    }
}
if (posisi == -1) cout << "tidak ada" << endl;
else cout << "indeks " << posisi << endl;
\end{Kode}
\end{column}
\begin{column}{0.41\textwidth}
\only<1>{\begin{TebakDulu}
{\ITERAKickerFont\fontsize{10pt}{12pt}\selectfont\color{ITERAAlert}TEBAK DULU}\par\vspace{4pt}
Sebelum dijalankan, tulis keluarannya di kertas.
\end{TebakDulu}
}\only<2>{\KeluaranFile[]{keluaran/P11/k001.txt}
\begin{Catatan}
Nilai awal -1 menandai belum ditemukan, karena -1 tidak pernah menjadi indeks yang sah.
\end{Catatan}
}
\end{column}
\end{columns}
\note<1>{Minta mahasiswa menebak keluaran sebelum membuka slide berikutnya. Tanyakan satu atau dua tebakan yang berbeda, lalu buktikan.}
\end{frame}

\begin{frame}[fragile]{Mencari semua kemunculan}
\framesubtitle{Tanpa break}
\begin{columns}[T]
\begin{column}{0.60\textwidth}
\begin{Kode}[listing options={style=ITERACodeKecil}]{semua.cpp}
int a[8] = {4, 9, 2, 7, 5, 9, 1, 8};
int x = 9, banyak = 0;
for (int i = 0; i < 8; i++) {
    if (a[i] == x) {
        cout << "indeks " << i << endl;
        banyak++;
    }
}
cout << "muncul " << banyak << " kali" << endl;
\end{Kode}
\end{column}
\begin{column}{0.36\textwidth}
\only<1>{\begin{TebakDulu}
{\ITERAKickerFont\fontsize{10pt}{12pt}\selectfont\color{ITERAAlert}TEBAK DULU}\par\vspace{4pt}
Sebelum dijalankan, tulis keluarannya di kertas.
\end{TebakDulu}
}\only<2>{\KeluaranFile[listing options={style=ITERAPlainKecil}]{keluaran/P11/k002.txt}
\begin{Catatan}
Untuk semua kemunculan, perulangan tidak dihentikan; pencacah mencatat banyaknya.
\end{Catatan}
}
\end{column}
\end{columns}
\note<1>{Minta mahasiswa menebak keluaran sebelum membuka slide berikutnya. Tanyakan satu atau dua tebakan yang berbeda, lalu buktikan.}
\end{frame}

\begin{frame}{Ide binary search}
\framesubtitle{Buang separuh setiap langkah}
Cari 21 pada \{3, 8, 15, 21, 34, 42, 57, 66, 79, 90\}, indeks 0 sampai 9.
\vspace{8pt}

\small
\begin{tabular}{@{}>{\raggedright\arraybackslash}p{110pt}>{\raggedright\arraybackslash}p{82pt}>{\raggedright\arraybackslash}p{82pt}>{\raggedright\arraybackslash}p{82pt}>{\raggedright\arraybackslash}p{110pt}>{\raggedright\arraybackslash}p{358pt}@{}}
\kepala{Langkah} & \kepala{low} & \kepala{high} & \kepala{mid} & \kepala{a[mid]} & \kepala{Keputusan}\\
1 & 0 & 9 & 4 & 34 & 34 > 21, cari di kiri: high = 3\\
\barisgenap 2 & 0 & 3 & 1 & 8 & 8 < 21, cari di kanan: low = 2\\
3 & 2 & 3 & 2 & 15 & 15 < 21, low = 3\\
\barisgenap 4 & 3 & 3 & 3 & 21 & ditemukan di indeks 3\\
\garisdasar
\end{tabular}
\vspace{10pt}

Binary search hanya benar pada data yang sudah terurut.
\end{frame}

\begin{frame}[fragile]{Binary search di C++}
\framesubtitle{low, high, mid}
\begin{columns}[T]
\begin{column}{0.60\textwidth}
\begin{Kode}[listing options={style=ITERACodeKecil}]{biner.cpp}
int a[10] = {3, 8, 15, 21, 34, 42, 57, 66, 79, 90};
int x = 21, low = 0, high = 9, posisi = -1;
while (low <= high) {
    int mid = (low + high) / 2;
    if (a[mid] == x) { posisi = mid; break; }
    else if (a[mid] < x) low = mid + 1;
    else high = mid - 1;
}
cout << "indeks " << posisi << endl;
\end{Kode}
\end{column}
\begin{column}{0.36\textwidth}
\only<1>{\begin{TebakDulu}
{\ITERAKickerFont\fontsize{10pt}{12pt}\selectfont\color{ITERAAlert}TEBAK DULU}\par\vspace{4pt}
Sebelum dijalankan, tulis keluarannya di kertas.
\end{TebakDulu}
}\only<2>{\KeluaranFile[listing options={style=ITERAPlainKecil}]{keluaran/P11/k003.txt}
\begin{Catatan}
Syaratnya \texttt{low <= high}. Dengan \texttt{<}, elemen terakhir yang tersisa tidak sempat diperiksa.
\end{Catatan}
}
\end{column}
\end{columns}
\note<1>{Minta mahasiswa menebak keluaran sebelum membuka slide berikutnya. Tanyakan satu atau dua tebakan yang berbeda, lalu buktikan.}
\end{frame}

\begin{frame}{Berapa langkahnya?}
\framesubtitle{Membandingkan dua algoritma}
\small
\begin{tabular}{@{}>{\raggedright\arraybackslash}p{249pt}>{\raggedright\arraybackslash}p{307pt}>{\raggedright\arraybackslash}p{307pt}@{}}
\kepala{Banyak data} & \kepala{Linear, kasus terburuk} & \kepala{Binary, kasus terburuk}\\
10 & 10 & 4\\
\barisgenap 100 & 100 & 7\\
1.000 & 1.000 & 10\\
\barisgenap 1.000.000 & 1.000.000 & 20\\
\garisdasar
\end{tabular}
\vspace{10pt}

Binary search jauh lebih cepat, tetapi butuh data terurut. Mengurutkan data adalah materi minggu depan.
\end{frame}

\begin{frame}[fragile]{Tebak keluarannya}
\framesubtitle{Latihan menjelaskan, 3 menit}
\begin{columns}[T]
\begin{column}{0.56\textwidth}
\begin{Kode}[listing options={style=ITERACodeKecil}]{tebak.cpp}
int a[7] = {2, 5, 8, 12, 16, 23, 38};
int x = 23, low = 0, high = 6, langkah = 0;
while (low <= high) {
    int mid = (low + high) / 2;
    langkah++;
    if (a[mid] == x) break;
    if (a[mid] < x) low = mid + 1;
    else high = mid - 1;
}
cout << langkah << " " << low << " " << high << endl;
\end{Kode}
\end{column}
\begin{column}{0.40\textwidth}
\begin{Catatan}
Tulis keluarannya di kertas, karakter demi karakter. Jangan dijalankan dulu.
\end{Catatan}
\end{column}
\end{columns}
\note{Contoh soal Bagian B. Beri waktu 3 menit sebelum slide jawaban.}
\end{frame}

\begin{frame}[fragile]{Tebak keluarannya: jawaban}
\framesubtitle{Latihan menjelaskan}
\begin{columns}[T]
\begin{column}{0.56\textwidth}
\begin{Kode}[listing options={style=ITERACodeKecil}]{tebak.cpp}
int a[7] = {2, 5, 8, 12, 16, 23, 38};
int x = 23, low = 0, high = 6, langkah = 0;
while (low <= high) {
    int mid = (low + high) / 2;
    langkah++;
    if (a[mid] == x) break;
    if (a[mid] < x) low = mid + 1;
    else high = mid - 1;
}
cout << langkah << " " << low << " " << high << endl;
\end{Kode}
\end{column}
\begin{column}{0.40\textwidth}
\begin{Keluaran}[]
2 4 6
\end{Keluaran}
\begin{Catatan}
mid = 3 (12 < 23, low = 4), lalu mid = 5 (23, ketemu). Dua langkah; saat berhenti low = 4 dan high = 6.
\end{Catatan}
\end{column}
\end{columns}
\note{Tanyakan jawaban yang paling sering salah, lalu telusuri bersama.}
\end{frame}

\begin{frame}{Error yang sering muncul minggu ini}
\framesubtitle{Pesan diambil langsung dari g++}
\footnotesize
\begin{tabular}{@{}>{\raggedright\arraybackslash}p{208pt}>{\raggedright\arraybackslash}p{256pt}>{\raggedright\arraybackslash}p{398pt}@{}}
\kepala{Kesalahan} & \kepala{Contoh} & \kepala{Pesan pertama dari g++}\\
Membandingkan array dengan == & \texttt{if (a == x)} & \texttt{error: ISO C++ forbids comparison between pointer and integer}\\
\barisgenap Variabel di luar jangkauan & \texttt{mid dipakai di luar while} & \texttt{error: 'mid' was not declared in this scope}\\
Kutip tunggal untuk string & \texttt{nama[i] == 'Rani'} & \texttt{error: no match for 'operator=='}\\
\barisgenap Posisi tidak dipakai & \texttt{int posisi = -1; tak dipakai} & \texttt{warning: unused variable 'posisi'}\\
Kurang titik koma di dalam if & \texttt{posisi = mid break;} & \texttt{error: expected ';' before 'break'}\\
\garisdasar
\end{tabular}
\vspace{10pt}

{\small Pesan di tabel ini dihasilkan dengan g++ dan opsi -Wall, sama seperti di cek.sh.}
\note{Minta mahasiswa memotret slide ini.}
\end{frame}

\Bagian{4}{Hands-on bertingkat}{45 menit, dari dasar sampai tantangan}
\note{Mahasiswa membuka folder 02 Hands-on/P11 - Searching - Hands-on.}

\begin{frame}{Peta hands-on}
\framesubtitle{Folder 02 Hands-on/P11 - Searching - Hands-on}
\small
\begin{tabular}{@{}>{\raggedright\arraybackslash}p{36pt}>{\raggedright\arraybackslash}p{267pt}>{\raggedright\arraybackslash}p{110pt}>{\raggedright\arraybackslash}p{83pt}>{\raggedright\arraybackslash}p{341pt}@{}}
\kepala{No} & \kepala{File} & \kepala{Tingkat} & \kepala{Waktu} & \kepala{Yang dilatih}\\
1 & \texttt{handson1-cari.cpp} & Dasar & 10' & linear search, hasil -1, hitung perbandingan\\
\barisgenap 2 & \texttt{handson2-nama.cpp} & Menengah & 15' & mencari semua kemunculan string\\
3 & \texttt{handson3-biner.cpp} & Tantangan & 20' & trace low-high-mid, perbaiki tiga kesalahan\\
\barisgenap + & \texttt{bonus-banding.cpp} & Pengayaan & sisa & menghitung perbandingan dua algoritma\\
\garisdasar
\end{tabular}
\vspace{10pt}

Kerjakan berurutan. Cocokkan keluaran dengan folder \texttt{expected/}; latihan yang membaca masukan memakai data di folder \texttt{input/}.\note{Saat berkeliling, minta mahasiswa yang mengerjakan latihan tantangan menjelaskan blok PENJELASAN mereka.}
\end{frame}

\begin{frame}{Cara mengecek dan aturannya}
\framesubtitle{Hands-on}
\begin{columns}[T]
\begin{column}{0.47\textwidth}
\begin{Blok}{Di GDB Online}
Salin isi file latihan, lengkapi TODO, klik Run. Bila program membaca masukan, ketik isi file di \texttt{input/}. Bandingkan dengan \texttt{expected/}.
\end{Blok}
\end{column}
\begin{column}{0.47\textwidth}
\begin{BlokGelap}{Di komputer sendiri}
Jalankan \texttt{bash cek.sh}. Skrip mengompilasi dengan \texttt{-Wall -Werror}, memberi masukan dari \texttt{input/}, lalu mencocokkan keluaran.
\end{BlokGelap}
\end{column}
\end{columns}
\vspace{6pt}
\begin{Catatan}
AI boleh untuk memahami pesan error, tidak untuk meminta solusi utuh. Exit Ticket dikerjakan tanpa bantuan apa pun.
\end{Catatan}
\note{Hands-on tidak dinilai langsung; yang dinilai Exit Ticket dan tugas.}
\end{frame}

\Bagian{5}{Exit Ticket dan tugas}{25 menit terakhir, lalu pekerjaan rumah}
\note{Mulai Exit Ticket tepat waktu.}

\begin{frame}{Exit Ticket 11}
\framesubtitle{Individual, 25 menit, tanpa AI dan internet}
\small
\begin{tabular}{@{}>{\raggedright\arraybackslash}p{60pt}>{\raggedright\arraybackslash}p{560pt}>{\raggedright\arraybackslash}p{80pt}>{\raggedright\arraybackslash}p{150pt}@{}}
\kepala{Soal} & \kepala{Isi} & \kepala{Poin} & \kepala{CPMK}\\
A1 & Tulis linear search yang menampilkan semua indeks kemunculan sebuah nilai & 25 & CPMK0607\\
\barisgenap A2 & Tulis binary search pada array terurut yang menampilkan indeks atau tidak ditemukan & 25 & CPMK0607\\
B1 & Isi trace table low, high, mid untuk sebuah pencarian & 20 & CPMK0608\\
\barisgenap B2 & Jelaskan mengapa binary search gagal pada data yang belum terurut & 20 & CPMK0608\\
B3 & Bandingkan banyak perbandingan dua algoritma untuk data yang diberikan & 10 & CPMK0608\\
\garisdasar
\end{tabular}
\vspace{10pt}

{\small Soal A dinilai dari keluaran yang tepat sama. Soal B dinilai dari ketepatan dan alasan. Pelanggaran aturan membuat nilai Exit Ticket ini 0.}
\note{Soal lengkap dibagikan terpisah dan tidak disimpan di repo publik.}
\end{frame}

\begin{frame}{Tugas 11: Direktori Kontak}
\framesubtitle{Dikumpulkan sebelum Pertemuan 12}
\begin{columns}[T]
\begin{column}{0.47\textwidth}
\begin{Blok}{Bagian A, program, 50 poin}
Buat program direktori kontak dengan menu berikut. Data dimasukkan dalam urutan abjad nama; program menolak nama yang tidak urut. Ketentuannya di slide berikut.
\end{Blok}
\end{column}
\begin{column}{0.47\textwidth}
\begin{BlokGelap}{Bagian B, penjelasan, 50 poin}
Komentar maksimal 150 kata: trace table binary search untuk satu nama yang ada dan satu yang tidak ada, serta perbandingan banyak langkah linear dan binary search pada data kalian. Jelaskan mengapa data harus urut abjad.
\end{BlokGelap}
\end{column}
\end{columns}
\vspace{8pt}
{\small Data di tugas adalah data kalian sendiri. Rubrik lengkap ada di Modul Belajar 11.}
\note{Ingatkan bahwa penjelasan disalin dari AI mudah dikenali karena tidak menyebut pengalaman mereka sendiri.}
\end{frame}

\begin{frame}{Ketentuan Tugas 11}
\framesubtitle{Direktori Kontak}
\small
\begin{tabular}{@{}>{\raggedright\arraybackslash}p{87pt}>{\raggedright\arraybackslash}p{788pt}@{}}
\kepala{No} & \kepala{Ketentuan Bagian A}\\
1 & Tambah kontak, maksimal 10 (nama satu kata dan nomor telepon), dengan validasi urutan abjad\\
\barisgenap 2 & Tampilkan semua kontak\\
3 & Cari nama dengan linear search; tampilkan nomor dan banyak perbandingan\\
\barisgenap 4 & Cari nama dengan binary search; tampilkan nomor dan banyak perbandingan\\
5 & Keluar\\
\barisgenap Bonus & Tambahkan pencarian awalan nama, misalnya semua kontak yang diawali "Ra".\\
\garisdasar
\end{tabular}
\note{Bacakan ketentuan satu per satu dan tanyakan apakah ada yang belum jelas.}
\end{frame}

\begin{frame}{Rubrik Tugas 11}
\framesubtitle{Empat level, sama di semua instrumen}
\footnotesize
\begin{tabular}{@{}>{\raggedright\arraybackslash}p{166pt}>{\raggedright\arraybackslash}p{44pt}>{\raggedright\arraybackslash}p{166pt}>{\raggedright\arraybackslash}p{156pt}>{\raggedright\arraybackslash}p{146pt}>{\raggedright\arraybackslash}p{146pt}@{}}
\kepala{Kriteria} & \kepala{Poin} & \kepala{Sangat baik 85--100} & \kepala{Baik 70--84} & \kepala{Cukup 55--69} & \kepala{Kurang <55}\\
A1 Program berjalan & 20 & tanpa error dan peringatan & ada peringatan & perlu satu perbaikan & tidak terkompilasi\\
\barisgenap A2 Kelengkapan menu & 15 & lima menu dan validasi urutan & satu menu kurang & dua kurang & tanpa binary search\\
A3 Ketepatan pencarian & 15 & kedua pencarian tepat, termasuk tidak ditemukan & satu kasus salah & dua salah & pencarian salah\\
\barisgenap B1 Trace table dan perbandingan & 30 & dua trace tepat dan analisis jelas & satu trace keliru & tanpa analisis & tidak ada atau disalin\\
B2 Cerita error dan pengujian & 20 & pesan, sebab, perbaikan, uji & pesan tidak dikutip & hanya perbaikan & tidak ada\\
\garisdasar
\end{tabular}
\vspace{8pt}

{\small Nilai kriteria = poin × skor level. Bagian A total 50, Bagian B total 50.}
\note{Rubrik ini sama strukturnya setiap minggu; yang berubah hanya isi kriteria A2, A3, dan B1.}
\end{frame}

\begin{frame}{Sebelum Pertemuan 12}
\framesubtitle{Persiapan}
\begin{columns}[T]
\begin{column}{0.31\textwidth}
\begin{Langkah}{1}{Selesaikan}
Hands-on yang belum lulus \texttt{cek.sh}, terutama latihan tantangan.
\end{Langkah}
\end{column}
\begin{column}{0.31\textwidth}
\begin{Langkah}{2}{Kumpulkan}
Tugas 11 sebelum Pertemuan 12 dimulai.
\end{Langkah}
\end{column}
\begin{column}{0.31\textwidth}
\begin{Langkah}{3}{Baca}
Modul Belajar 11 untuk mengulang, lalu bagian awal modul berikutnya.
\end{Langkah}
\end{column}
\end{columns}
\vspace{10pt}
Pertemuan 12 membahas sorting: bubble sort dan selection sort, supaya data bisa diurutkan sebelum dicari dengan binary search.\note{Tanyakan satu hal yang paling membingungkan hari ini untuk dibuka di review minggu depan.}
\end{frame}

\SlidePenutup{Terima kasih}{Sampai jumpa di Pertemuan 12: sorting}
\note{Slide penutup.}
\end{document}
"
\documentclass[t]{beamer}
\usetheme{ITERA}
% Mode catatan pembicara: aktif bila \catatan didefinisikan sebelum \input.
\ifdefined\catatan
  \setbeameroption{show only notes}
  \setbeamertemplate{note page}{\vspace*{20pt}\hspace*{4pt}\begin{minipage}{900pt}
    {\ITERAKickerFont\fontsize{11pt}{14pt}\selectfont\color{ITERAGoldText}CATATAN PEMBICARA \ \ |\ \ SLIDE \insertframenumber}\par\vspace{4pt}
    {\ITERADisplay\bfseries\fontsize{22pt}{28pt}\selectfont\color{ITERAStructure}\insertframetitle}\par\vspace{12pt}
    \fontsize{15pt}{23pt}\selectfont\insertnote\end{minipage}}
\fi
\tikzset{fase/.style={draw=none,fill=#1,rounded corners=3pt,minimum width=158pt,minimum height=118pt,text width=146pt,align=center}}

\renewcommand\ITERAmeeting{Pertemuan 11}
\begin{document}
\SlideJudul{IF25-11002 PRAKTIKUM PEMROGRAMAN}{Pertemuan 11\\Searching}{Mencari data di array dengan linear search dan binary search, dan membandingkan banyak langkahnya}{Tim Pengajar}{Tahun 2026. Sejajar dengan Pertemuan 11 Algoritma Pemrograman: Searching.}
\note{Buka dengan menanyakan satu hal dari pertemuan lalu yang masih membingungkan.}

\begin{frame}{Dua mata kuliah, satu minggu}
\framesubtitle{Sebelum mulai}
Topik minggu ini sama di dua mata kuliah, dengan peran yang berbeda.
\vspace{10pt}

\begin{columns}[T]
\begin{column}{0.47\textwidth}
\begin{Blok}{Algoritma: Searching}
Algoritma pencarian beruntun dan biner, prasyarat data terurut, serta menghitung banyak perbandingan.
\end{Blok}
\end{column}
\begin{column}{0.47\textwidth}
\begin{BlokGelap}{Praktikum: Searching}
Menulis linear search dan binary search di C++, menelusuri low, high, dan mid, serta menghitung perbandingan.
\end{BlokGelap}
\end{column}
\end{columns}
\note{Tanyakan apa yang sudah dibahas di Algoritma minggu ini. Gunakan jawabannya sebagai jembatan ke praktikum.}
\end{frame}

\begin{frame}{Saat pulang nanti, kalian bisa}
\framesubtitle{Target hari ini}
\begin{columns}[T]
\begin{column}{0.47\textwidth}
\begin{Langkah}{1}{Menerapkan}
Menulis program yang mencari data pada array dengan linear search dan binary search, termasuk mencari semua kemunculan dan menangani data yang tidak ditemukan

{\footnotesize\color{ITERAInkMuted}Sub-CPMK 11.1, CPMK0607}
\end{Langkah}
\end{column}
\begin{column}{0.47\textwidth}
\begin{Langkah}{2}{Menjelaskan}
Menelusuri langkah binary search dengan trace table low, high, dan mid, serta menjelaskan perbedaan banyak perbandingan kedua algoritma

{\footnotesize\color{ITERAInkMuted}Sub-CPMK 11.2, CPMK0608}
\end{Langkah}
\end{column}
\end{columns}
\vspace{8pt}
\begin{Catatan}
Dua kemampuan ini dinilai sama berat di Exit Ticket, tugas, UTS, dan UAS.
\end{Catatan}
\note{Bacakan target dengan pelan. Kata kuncinya menerapkan dan menjelaskan.}
\end{frame}

\begin{frame}{Ingat minggu lalu}
\framesubtitle{Review, 5 menit}
\begin{columns}[T]
\begin{column}{0.31\textwidth}
\begin{BlokGelap}{Pertanyaan 1}
Bagaimana menghitung rata-rata setiap kolom array 2D?
\end{BlokGelap}
\end{column}
\begin{column}{0.31\textwidth}
\begin{BlokGelap}{Pertanyaan 2}
Apa elemen diagonal utama matriks?
\end{BlokGelap}
\end{column}
\begin{column}{0.31\textwidth}
\begin{BlokGelap}{Pertanyaan 3}
Apa hasil \texttt{s.find("tik")} untuk \texttt{"Praktikum"}?
\end{BlokGelap}
\end{column}
\end{columns}
\vspace{8pt}
Jawab di kertas dulu, lalu cocokkan dengan teman sebelah.\note{Pilih dua mahasiswa untuk menjawab. Koreksi singkat, jangan mengajar ulang.}
\end{frame}

\Bagian{1}{Masalah pembuka}{Mencari nama di daftar wisuda}
\note{Masalah pembuka 10 menit.}

\begin{frame}[fragile]{Mencari nama di daftar wisuda}
\framesubtitle{Masalah minggu ini}
\begin{columns}[T]
\begin{column}{0.55\textwidth}
{Panitia wisuda memegang daftar 1.200 nama yang sudah urut abjad. Seorang orang tua bertanya, apakah anaknya ada di daftar?

\vspace{6pt}
Membaca dari nama pertama sampai ketemu bisa butuh ratusan kali melihat. Padahal saat mencari kata di kamus, kita tidak membaca dari halaman pertama.\par}
\vspace{6pt}
\begin{Catatan}
\textbf{Diskusi 2 menit.} Bagaimana kalian mencari kata di kamus? Mengapa cara itu tidak bisa dipakai pada daftar yang belum urut?
\end{Catatan}
\end{column}
\begin{column}{0.41\textwidth}
\begin{Keluaran}[]
Cari: Siti Rahma
Ditemukan di nomor 1023
Linear search : 1023 perbandingan
Binary search : 10 perbandingan
\end{Keluaran}
\end{column}
\end{columns}
\note{Tampung beberapa jawaban tanpa menilai benar salah.}
\end{frame}

\begin{frame}{Dari langkah ke instruksi}
\framesubtitle{Sambungan dengan Algoritma}
\begin{columns}[T]
\begin{column}{0.31\textwidth}
\begin{Langkah}{1}{Bandingkan satu per satu}
Cara paling sederhana: periksa setiap elemen dari depan.
\end{Langkah}
\end{column}
\begin{column}{0.31\textwidth}
\begin{Langkah}{2}{Manfaatkan urutan}
Bila data urut, buang separuh daftar di setiap langkah.
\end{Langkah}
\end{column}
\begin{column}{0.31\textwidth}
\begin{Langkah}{3}{Terjemahkan}
Linear search dengan \texttt{for}, binary search dengan \texttt{while} dan tiga indeks.
\end{Langkah}
\end{column}
\end{columns}
\vspace{10pt}
Langkah 1 dan 2 dilatih di Algoritma. Langkah 3 kita kerjakan hari ini dengan C++.\note{Hubungkan eksplisit ke Algoritma minggu ini.}
\end{frame}

\Bagian{2}{Coba dulu, baru dijelaskan}{25 menit eksperimen di GDB Online}
\note{Struggle 25 menit. Berkeliling, jangan menjelaskan materi dulu.}

\begin{frame}{Misi kalian}
\framesubtitle{Struggle, 25 menit}
\begin{columns}[T]
\begin{column}{0.31\textwidth}
\begin{Langkah}{1}{Cari}
Diberikan \texttt{int a[8] = \{4, 9, 2, 7, 5, 9, 1, 8\};}, tampilkan indeks angka 7.
\end{Langkah}
\end{column}
\begin{column}{0.31\textwidth}
\begin{Langkah}{2}{Semua}
Tampilkan semua indeks angka 9, bukan hanya yang pertama.
\end{Langkah}
\end{column}
\begin{column}{0.31\textwidth}
\begin{Langkah}{3}{Tidak ada}
Apa yang ditampilkan program kalian bila mencari angka 6?
\end{Langkah}
\end{column}
\end{columns}
\vspace{8pt}
\begin{columns}[T]
\begin{column}{0.47\textwidth}
\begin{Blok}{Boleh}
Bertanya ke teman, mengubah apa saja, sengaja membuat error.
\end{Blok}
\end{column}
\begin{column}{0.47\textwidth}
\begin{BlokAlert}{Belum dulu}
Mencari jawaban jadi di internet atau AI.
\end{BlokAlert}
\end{column}
\end{columns}
\note{Catat kesalahan yang sering muncul untuk dibahas di debrief.}
\end{frame}

\begin{frame}{Apa yang kalian temukan?}
\framesubtitle{Debrief struggle}
\begin{columns}[T]
\begin{column}{0.31\textwidth}
\begin{BlokGelap}{Pertanyaan 1}
Bagaimana program tahu sebuah data tidak ditemukan?
\end{BlokGelap}
\end{column}
\begin{column}{0.31\textwidth}
\begin{BlokGelap}{Pertanyaan 2}
Kapan \texttt{break} dipakai dan kapan tidak?
\end{BlokGelap}
\end{column}
\begin{column}{0.31\textwidth}
\begin{BlokGelap}{Pertanyaan 3}
Berapa perbandingan untuk mencari elemen terakhir?
\end{BlokGelap}
\end{column}
\end{columns}
\vspace{10pt}
Jawaban kalian akan kita uji di bagian berikutnya.\note{Tulis jawaban di papan; buktikan atau patahkan di bagian materi.}
\end{frame}

\Bagian{3}{Mencari dengan cerdas}{Linear search, binary search, dan menghitung langkah}
\note{Materi 40 menit. Tunjukkan setiap contoh langsung di GDB Online.}

\begin{frame}[fragile]{Linear search}
\framesubtitle{Periksa satu per satu}
\begin{columns}[T]
\begin{column}{0.55\textwidth}
\begin{Kode}[]{linear.cpp}
int a[8] = {4, 9, 2, 7, 5, 9, 1, 8};
int x = 7, posisi = -1;
for (int i = 0; i < 8; i++) {
    if (a[i] == x) {
        posisi = i;
        break;
    }
}
if (posisi == -1) cout << "tidak ada" << endl;
else cout << "indeks " << posisi << endl;
\end{Kode}
\end{column}
\begin{column}{0.41\textwidth}
\only<1>{\begin{TebakDulu}
{\ITERAKickerFont\fontsize{10pt}{12pt}\selectfont\color{ITERAAlert}TEBAK DULU}\par\vspace{4pt}
Sebelum dijalankan, tulis keluarannya di kertas.
\end{TebakDulu}
}\only<2>{\KeluaranFile[]{keluaran/P11/k001.txt}
\begin{Catatan}
Nilai awal -1 menandai belum ditemukan, karena -1 tidak pernah menjadi indeks yang sah.
\end{Catatan}
}
\end{column}
\end{columns}
\note<1>{Minta mahasiswa menebak keluaran sebelum membuka slide berikutnya. Tanyakan satu atau dua tebakan yang berbeda, lalu buktikan.}
\end{frame}

\begin{frame}[fragile]{Mencari semua kemunculan}
\framesubtitle{Tanpa break}
\begin{columns}[T]
\begin{column}{0.60\textwidth}
\begin{Kode}[listing options={style=ITERACodeKecil}]{semua.cpp}
int a[8] = {4, 9, 2, 7, 5, 9, 1, 8};
int x = 9, banyak = 0;
for (int i = 0; i < 8; i++) {
    if (a[i] == x) {
        cout << "indeks " << i << endl;
        banyak++;
    }
}
cout << "muncul " << banyak << " kali" << endl;
\end{Kode}
\end{column}
\begin{column}{0.36\textwidth}
\only<1>{\begin{TebakDulu}
{\ITERAKickerFont\fontsize{10pt}{12pt}\selectfont\color{ITERAAlert}TEBAK DULU}\par\vspace{4pt}
Sebelum dijalankan, tulis keluarannya di kertas.
\end{TebakDulu}
}\only<2>{\KeluaranFile[listing options={style=ITERAPlainKecil}]{keluaran/P11/k002.txt}
\begin{Catatan}
Untuk semua kemunculan, perulangan tidak dihentikan; pencacah mencatat banyaknya.
\end{Catatan}
}
\end{column}
\end{columns}
\note<1>{Minta mahasiswa menebak keluaran sebelum membuka slide berikutnya. Tanyakan satu atau dua tebakan yang berbeda, lalu buktikan.}
\end{frame}

\begin{frame}{Ide binary search}
\framesubtitle{Buang separuh setiap langkah}
Cari 21 pada \{3, 8, 15, 21, 34, 42, 57, 66, 79, 90\}, indeks 0 sampai 9.
\vspace{8pt}

\small
\begin{tabular}{@{}>{\raggedright\arraybackslash}p{110pt}>{\raggedright\arraybackslash}p{82pt}>{\raggedright\arraybackslash}p{82pt}>{\raggedright\arraybackslash}p{82pt}>{\raggedright\arraybackslash}p{110pt}>{\raggedright\arraybackslash}p{358pt}@{}}
\kepala{Langkah} & \kepala{low} & \kepala{high} & \kepala{mid} & \kepala{a[mid]} & \kepala{Keputusan}\\
1 & 0 & 9 & 4 & 34 & 34 > 21, cari di kiri: high = 3\\
\barisgenap 2 & 0 & 3 & 1 & 8 & 8 < 21, cari di kanan: low = 2\\
3 & 2 & 3 & 2 & 15 & 15 < 21, low = 3\\
\barisgenap 4 & 3 & 3 & 3 & 21 & ditemukan di indeks 3\\
\garisdasar
\end{tabular}
\vspace{10pt}

Binary search hanya benar pada data yang sudah terurut.
\end{frame}

\begin{frame}[fragile]{Binary search di C++}
\framesubtitle{low, high, mid}
\begin{columns}[T]
\begin{column}{0.60\textwidth}
\begin{Kode}[listing options={style=ITERACodeKecil}]{biner.cpp}
int a[10] = {3, 8, 15, 21, 34, 42, 57, 66, 79, 90};
int x = 21, low = 0, high = 9, posisi = -1;
while (low <= high) {
    int mid = (low + high) / 2;
    if (a[mid] == x) { posisi = mid; break; }
    else if (a[mid] < x) low = mid + 1;
    else high = mid - 1;
}
cout << "indeks " << posisi << endl;
\end{Kode}
\end{column}
\begin{column}{0.36\textwidth}
\only<1>{\begin{TebakDulu}
{\ITERAKickerFont\fontsize{10pt}{12pt}\selectfont\color{ITERAAlert}TEBAK DULU}\par\vspace{4pt}
Sebelum dijalankan, tulis keluarannya di kertas.
\end{TebakDulu}
}\only<2>{\KeluaranFile[listing options={style=ITERAPlainKecil}]{keluaran/P11/k003.txt}
\begin{Catatan}
Syaratnya \texttt{low <= high}. Dengan \texttt{<}, elemen terakhir yang tersisa tidak sempat diperiksa.
\end{Catatan}
}
\end{column}
\end{columns}
\note<1>{Minta mahasiswa menebak keluaran sebelum membuka slide berikutnya. Tanyakan satu atau dua tebakan yang berbeda, lalu buktikan.}
\end{frame}

\begin{frame}{Berapa langkahnya?}
\framesubtitle{Membandingkan dua algoritma}
\small
\begin{tabular}{@{}>{\raggedright\arraybackslash}p{249pt}>{\raggedright\arraybackslash}p{307pt}>{\raggedright\arraybackslash}p{307pt}@{}}
\kepala{Banyak data} & \kepala{Linear, kasus terburuk} & \kepala{Binary, kasus terburuk}\\
10 & 10 & 4\\
\barisgenap 100 & 100 & 7\\
1.000 & 1.000 & 10\\
\barisgenap 1.000.000 & 1.000.000 & 20\\
\garisdasar
\end{tabular}
\vspace{10pt}

Binary search jauh lebih cepat, tetapi butuh data terurut. Mengurutkan data adalah materi minggu depan.
\end{frame}

\begin{frame}[fragile]{Tebak keluarannya}
\framesubtitle{Latihan menjelaskan, 3 menit}
\begin{columns}[T]
\begin{column}{0.56\textwidth}
\begin{Kode}[listing options={style=ITERACodeKecil}]{tebak.cpp}
int a[7] = {2, 5, 8, 12, 16, 23, 38};
int x = 23, low = 0, high = 6, langkah = 0;
while (low <= high) {
    int mid = (low + high) / 2;
    langkah++;
    if (a[mid] == x) break;
    if (a[mid] < x) low = mid + 1;
    else high = mid - 1;
}
cout << langkah << " " << low << " " << high << endl;
\end{Kode}
\end{column}
\begin{column}{0.40\textwidth}
\begin{Catatan}
Tulis keluarannya di kertas, karakter demi karakter. Jangan dijalankan dulu.
\end{Catatan}
\end{column}
\end{columns}
\note{Contoh soal Bagian B. Beri waktu 3 menit sebelum slide jawaban.}
\end{frame}

\begin{frame}[fragile]{Tebak keluarannya: jawaban}
\framesubtitle{Latihan menjelaskan}
\begin{columns}[T]
\begin{column}{0.56\textwidth}
\begin{Kode}[listing options={style=ITERACodeKecil}]{tebak.cpp}
int a[7] = {2, 5, 8, 12, 16, 23, 38};
int x = 23, low = 0, high = 6, langkah = 0;
while (low <= high) {
    int mid = (low + high) / 2;
    langkah++;
    if (a[mid] == x) break;
    if (a[mid] < x) low = mid + 1;
    else high = mid - 1;
}
cout << langkah << " " << low << " " << high << endl;
\end{Kode}
\end{column}
\begin{column}{0.40\textwidth}
\begin{Keluaran}[]
2 4 6
\end{Keluaran}
\begin{Catatan}
mid = 3 (12 < 23, low = 4), lalu mid = 5 (23, ketemu). Dua langkah; saat berhenti low = 4 dan high = 6.
\end{Catatan}
\end{column}
\end{columns}
\note{Tanyakan jawaban yang paling sering salah, lalu telusuri bersama.}
\end{frame}

\begin{frame}{Error yang sering muncul minggu ini}
\framesubtitle{Pesan diambil langsung dari g++}
\footnotesize
\begin{tabular}{@{}>{\raggedright\arraybackslash}p{208pt}>{\raggedright\arraybackslash}p{256pt}>{\raggedright\arraybackslash}p{398pt}@{}}
\kepala{Kesalahan} & \kepala{Contoh} & \kepala{Pesan pertama dari g++}\\
Membandingkan array dengan == & \texttt{if (a == x)} & \texttt{error: ISO C++ forbids comparison between pointer and integer}\\
\barisgenap Variabel di luar jangkauan & \texttt{mid dipakai di luar while} & \texttt{error: 'mid' was not declared in this scope}\\
Kutip tunggal untuk string & \texttt{nama[i] == 'Rani'} & \texttt{error: no match for 'operator=='}\\
\barisgenap Posisi tidak dipakai & \texttt{int posisi = -1; tak dipakai} & \texttt{warning: unused variable 'posisi'}\\
Kurang titik koma di dalam if & \texttt{posisi = mid break;} & \texttt{error: expected ';' before 'break'}\\
\garisdasar
\end{tabular}
\vspace{10pt}

{\small Pesan di tabel ini dihasilkan dengan g++ dan opsi -Wall, sama seperti di cek.sh.}
\note{Minta mahasiswa memotret slide ini.}
\end{frame}

\Bagian{4}{Hands-on bertingkat}{45 menit, dari dasar sampai tantangan}
\note{Mahasiswa membuka folder 02 Hands-on/P11 - Searching - Hands-on.}

\begin{frame}{Peta hands-on}
\framesubtitle{Folder 02 Hands-on/P11 - Searching - Hands-on}
\small
\begin{tabular}{@{}>{\raggedright\arraybackslash}p{36pt}>{\raggedright\arraybackslash}p{267pt}>{\raggedright\arraybackslash}p{110pt}>{\raggedright\arraybackslash}p{83pt}>{\raggedright\arraybackslash}p{341pt}@{}}
\kepala{No} & \kepala{File} & \kepala{Tingkat} & \kepala{Waktu} & \kepala{Yang dilatih}\\
1 & \texttt{handson1-cari.cpp} & Dasar & 10' & linear search, hasil -1, hitung perbandingan\\
\barisgenap 2 & \texttt{handson2-nama.cpp} & Menengah & 15' & mencari semua kemunculan string\\
3 & \texttt{handson3-biner.cpp} & Tantangan & 20' & trace low-high-mid, perbaiki tiga kesalahan\\
\barisgenap + & \texttt{bonus-banding.cpp} & Pengayaan & sisa & menghitung perbandingan dua algoritma\\
\garisdasar
\end{tabular}
\vspace{10pt}

Kerjakan berurutan. Cocokkan keluaran dengan folder \texttt{expected/}; latihan yang membaca masukan memakai data di folder \texttt{input/}.\note{Saat berkeliling, minta mahasiswa yang mengerjakan latihan tantangan menjelaskan blok PENJELASAN mereka.}
\end{frame}

\begin{frame}{Cara mengecek dan aturannya}
\framesubtitle{Hands-on}
\begin{columns}[T]
\begin{column}{0.47\textwidth}
\begin{Blok}{Di GDB Online}
Salin isi file latihan, lengkapi TODO, klik Run. Bila program membaca masukan, ketik isi file di \texttt{input/}. Bandingkan dengan \texttt{expected/}.
\end{Blok}
\end{column}
\begin{column}{0.47\textwidth}
\begin{BlokGelap}{Di komputer sendiri}
Jalankan \texttt{bash cek.sh}. Skrip mengompilasi dengan \texttt{-Wall -Werror}, memberi masukan dari \texttt{input/}, lalu mencocokkan keluaran.
\end{BlokGelap}
\end{column}
\end{columns}
\vspace{6pt}
\begin{Catatan}
AI boleh untuk memahami pesan error, tidak untuk meminta solusi utuh. Exit Ticket dikerjakan tanpa bantuan apa pun.
\end{Catatan}
\note{Hands-on tidak dinilai langsung; yang dinilai Exit Ticket dan tugas.}
\end{frame}

\Bagian{5}{Exit Ticket dan tugas}{25 menit terakhir, lalu pekerjaan rumah}
\note{Mulai Exit Ticket tepat waktu.}

\begin{frame}{Exit Ticket 11}
\framesubtitle{Individual, 25 menit, tanpa AI dan internet}
\small
\begin{tabular}{@{}>{\raggedright\arraybackslash}p{60pt}>{\raggedright\arraybackslash}p{560pt}>{\raggedright\arraybackslash}p{80pt}>{\raggedright\arraybackslash}p{150pt}@{}}
\kepala{Soal} & \kepala{Isi} & \kepala{Poin} & \kepala{CPMK}\\
A1 & Tulis linear search yang menampilkan semua indeks kemunculan sebuah nilai & 25 & CPMK0607\\
\barisgenap A2 & Tulis binary search pada array terurut yang menampilkan indeks atau tidak ditemukan & 25 & CPMK0607\\
B1 & Isi trace table low, high, mid untuk sebuah pencarian & 20 & CPMK0608\\
\barisgenap B2 & Jelaskan mengapa binary search gagal pada data yang belum terurut & 20 & CPMK0608\\
B3 & Bandingkan banyak perbandingan dua algoritma untuk data yang diberikan & 10 & CPMK0608\\
\garisdasar
\end{tabular}
\vspace{10pt}

{\small Soal A dinilai dari keluaran yang tepat sama. Soal B dinilai dari ketepatan dan alasan. Pelanggaran aturan membuat nilai Exit Ticket ini 0.}
\note{Soal lengkap dibagikan terpisah dan tidak disimpan di repo publik.}
\end{frame}

\begin{frame}{Tugas 11: Direktori Kontak}
\framesubtitle{Dikumpulkan sebelum Pertemuan 12}
\begin{columns}[T]
\begin{column}{0.47\textwidth}
\begin{Blok}{Bagian A, program, 50 poin}
Buat program direktori kontak dengan menu berikut. Data dimasukkan dalam urutan abjad nama; program menolak nama yang tidak urut. Ketentuannya di slide berikut.
\end{Blok}
\end{column}
\begin{column}{0.47\textwidth}
\begin{BlokGelap}{Bagian B, penjelasan, 50 poin}
Komentar maksimal 150 kata: trace table binary search untuk satu nama yang ada dan satu yang tidak ada, serta perbandingan banyak langkah linear dan binary search pada data kalian. Jelaskan mengapa data harus urut abjad.
\end{BlokGelap}
\end{column}
\end{columns}
\vspace{8pt}
{\small Data di tugas adalah data kalian sendiri. Rubrik lengkap ada di Modul Belajar 11.}
\note{Ingatkan bahwa penjelasan disalin dari AI mudah dikenali karena tidak menyebut pengalaman mereka sendiri.}
\end{frame}

\begin{frame}{Ketentuan Tugas 11}
\framesubtitle{Direktori Kontak}
\small
\begin{tabular}{@{}>{\raggedright\arraybackslash}p{87pt}>{\raggedright\arraybackslash}p{788pt}@{}}
\kepala{No} & \kepala{Ketentuan Bagian A}\\
1 & Tambah kontak, maksimal 10 (nama satu kata dan nomor telepon), dengan validasi urutan abjad\\
\barisgenap 2 & Tampilkan semua kontak\\
3 & Cari nama dengan linear search; tampilkan nomor dan banyak perbandingan\\
\barisgenap 4 & Cari nama dengan binary search; tampilkan nomor dan banyak perbandingan\\
5 & Keluar\\
\barisgenap Bonus & Tambahkan pencarian awalan nama, misalnya semua kontak yang diawali "Ra".\\
\garisdasar
\end{tabular}
\note{Bacakan ketentuan satu per satu dan tanyakan apakah ada yang belum jelas.}
\end{frame}

\begin{frame}{Rubrik Tugas 11}
\framesubtitle{Empat level, sama di semua instrumen}
\footnotesize
\begin{tabular}{@{}>{\raggedright\arraybackslash}p{166pt}>{\raggedright\arraybackslash}p{44pt}>{\raggedright\arraybackslash}p{166pt}>{\raggedright\arraybackslash}p{156pt}>{\raggedright\arraybackslash}p{146pt}>{\raggedright\arraybackslash}p{146pt}@{}}
\kepala{Kriteria} & \kepala{Poin} & \kepala{Sangat baik 85--100} & \kepala{Baik 70--84} & \kepala{Cukup 55--69} & \kepala{Kurang <55}\\
A1 Program berjalan & 20 & tanpa error dan peringatan & ada peringatan & perlu satu perbaikan & tidak terkompilasi\\
\barisgenap A2 Kelengkapan menu & 15 & lima menu dan validasi urutan & satu menu kurang & dua kurang & tanpa binary search\\
A3 Ketepatan pencarian & 15 & kedua pencarian tepat, termasuk tidak ditemukan & satu kasus salah & dua salah & pencarian salah\\
\barisgenap B1 Trace table dan perbandingan & 30 & dua trace tepat dan analisis jelas & satu trace keliru & tanpa analisis & tidak ada atau disalin\\
B2 Cerita error dan pengujian & 20 & pesan, sebab, perbaikan, uji & pesan tidak dikutip & hanya perbaikan & tidak ada\\
\garisdasar
\end{tabular}
\vspace{8pt}

{\small Nilai kriteria = poin × skor level. Bagian A total 50, Bagian B total 50.}
\note{Rubrik ini sama strukturnya setiap minggu; yang berubah hanya isi kriteria A2, A3, dan B1.}
\end{frame}

\begin{frame}{Sebelum Pertemuan 12}
\framesubtitle{Persiapan}
\begin{columns}[T]
\begin{column}{0.31\textwidth}
\begin{Langkah}{1}{Selesaikan}
Hands-on yang belum lulus \texttt{cek.sh}, terutama latihan tantangan.
\end{Langkah}
\end{column}
\begin{column}{0.31\textwidth}
\begin{Langkah}{2}{Kumpulkan}
Tugas 11 sebelum Pertemuan 12 dimulai.
\end{Langkah}
\end{column}
\begin{column}{0.31\textwidth}
\begin{Langkah}{3}{Baca}
Modul Belajar 11 untuk mengulang, lalu bagian awal modul berikutnya.
\end{Langkah}
\end{column}
\end{columns}
\vspace{10pt}
Pertemuan 12 membahas sorting: bubble sort dan selection sort, supaya data bisa diurutkan sebelum dicari dengan binary search.\note{Tanyakan satu hal yang paling membingungkan hari ini untuk dibuka di review minggu depan.}
\end{frame}

\SlidePenutup{Terima kasih}{Sampai jumpa di Pertemuan 12: sorting}
\note{Slide penutup.}
\end{document}
"
\documentclass[t]{beamer}
\usetheme{ITERA}
% Mode catatan pembicara: aktif bila \catatan didefinisikan sebelum \input.
\ifdefined\catatan
  \setbeameroption{show only notes}
  \setbeamertemplate{note page}{\vspace*{20pt}\hspace*{4pt}\begin{minipage}{900pt}
    {\ITERAKickerFont\fontsize{11pt}{14pt}\selectfont\color{ITERAGoldText}CATATAN PEMBICARA \ \ |\ \ SLIDE \insertframenumber}\par\vspace{4pt}
    {\ITERADisplay\bfseries\fontsize{22pt}{28pt}\selectfont\color{ITERAStructure}\insertframetitle}\par\vspace{12pt}
    \fontsize{15pt}{23pt}\selectfont\insertnote\end{minipage}}
\fi
\tikzset{fase/.style={draw=none,fill=#1,rounded corners=3pt,minimum width=158pt,minimum height=118pt,text width=146pt,align=center}}

\renewcommand\ITERAmeeting{Pertemuan 11}
\begin{document}
\SlideJudul{IF25-11002 PRAKTIKUM PEMROGRAMAN}{Pertemuan 11\\Searching}{Mencari data di array dengan linear search dan binary search, dan membandingkan banyak langkahnya}{Tim Pengajar}{Tahun 2026. Sejajar dengan Pertemuan 11 Algoritma Pemrograman: Searching.}
\note{Buka dengan menanyakan satu hal dari pertemuan lalu yang masih membingungkan.}

\begin{frame}{Dua mata kuliah, satu minggu}
\framesubtitle{Sebelum mulai}
Topik minggu ini sama di dua mata kuliah, dengan peran yang berbeda.
\vspace{10pt}

\begin{columns}[T]
\begin{column}{0.47\textwidth}
\begin{Blok}{Algoritma: Searching}
Algoritma pencarian beruntun dan biner, prasyarat data terurut, serta menghitung banyak perbandingan.
\end{Blok}
\end{column}
\begin{column}{0.47\textwidth}
\begin{BlokGelap}{Praktikum: Searching}
Menulis linear search dan binary search di C++, menelusuri low, high, dan mid, serta menghitung perbandingan.
\end{BlokGelap}
\end{column}
\end{columns}
\note{Tanyakan apa yang sudah dibahas di Algoritma minggu ini. Gunakan jawabannya sebagai jembatan ke praktikum.}
\end{frame}

\begin{frame}{Saat pulang nanti, kalian bisa}
\framesubtitle{Target hari ini}
\begin{columns}[T]
\begin{column}{0.47\textwidth}
\begin{Langkah}{1}{Menerapkan}
Menulis program yang mencari data pada array dengan linear search dan binary search, termasuk mencari semua kemunculan dan menangani data yang tidak ditemukan

{\footnotesize\color{ITERAInkMuted}Sub-CPMK 11.1, CPMK0607}
\end{Langkah}
\end{column}
\begin{column}{0.47\textwidth}
\begin{Langkah}{2}{Menjelaskan}
Menelusuri langkah binary search dengan trace table low, high, dan mid, serta menjelaskan perbedaan banyak perbandingan kedua algoritma

{\footnotesize\color{ITERAInkMuted}Sub-CPMK 11.2, CPMK0608}
\end{Langkah}
\end{column}
\end{columns}
\vspace{8pt}
\begin{Catatan}
Dua kemampuan ini dinilai sama berat di Exit Ticket, tugas, UTS, dan UAS.
\end{Catatan}
\note{Bacakan target dengan pelan. Kata kuncinya menerapkan dan menjelaskan.}
\end{frame}

\begin{frame}{Ingat minggu lalu}
\framesubtitle{Review, 5 menit}
\begin{columns}[T]
\begin{column}{0.31\textwidth}
\begin{BlokGelap}{Pertanyaan 1}
Bagaimana menghitung rata-rata setiap kolom array 2D?
\end{BlokGelap}
\end{column}
\begin{column}{0.31\textwidth}
\begin{BlokGelap}{Pertanyaan 2}
Apa elemen diagonal utama matriks?
\end{BlokGelap}
\end{column}
\begin{column}{0.31\textwidth}
\begin{BlokGelap}{Pertanyaan 3}
Apa hasil \texttt{s.find("tik")} untuk \texttt{"Praktikum"}?
\end{BlokGelap}
\end{column}
\end{columns}
\vspace{8pt}
Jawab di kertas dulu, lalu cocokkan dengan teman sebelah.\note{Pilih dua mahasiswa untuk menjawab. Koreksi singkat, jangan mengajar ulang.}
\end{frame}

\Bagian{1}{Masalah pembuka}{Mencari nama di daftar wisuda}
\note{Masalah pembuka 10 menit.}

\begin{frame}[fragile]{Mencari nama di daftar wisuda}
\framesubtitle{Masalah minggu ini}
\begin{columns}[T]
\begin{column}{0.55\textwidth}
{Panitia wisuda memegang daftar 1.200 nama yang sudah urut abjad. Seorang orang tua bertanya, apakah anaknya ada di daftar?

\vspace{6pt}
Membaca dari nama pertama sampai ketemu bisa butuh ratusan kali melihat. Padahal saat mencari kata di kamus, kita tidak membaca dari halaman pertama.\par}
\vspace{6pt}
\begin{Catatan}
\textbf{Diskusi 2 menit.} Bagaimana kalian mencari kata di kamus? Mengapa cara itu tidak bisa dipakai pada daftar yang belum urut?
\end{Catatan}
\end{column}
\begin{column}{0.41\textwidth}
\begin{Keluaran}[]
Cari: Siti Rahma
Ditemukan di nomor 1023
Linear search : 1023 perbandingan
Binary search : 10 perbandingan
\end{Keluaran}
\end{column}
\end{columns}
\note{Tampung beberapa jawaban tanpa menilai benar salah.}
\end{frame}

\begin{frame}{Dari langkah ke instruksi}
\framesubtitle{Sambungan dengan Algoritma}
\begin{columns}[T]
\begin{column}{0.31\textwidth}
\begin{Langkah}{1}{Bandingkan satu per satu}
Cara paling sederhana: periksa setiap elemen dari depan.
\end{Langkah}
\end{column}
\begin{column}{0.31\textwidth}
\begin{Langkah}{2}{Manfaatkan urutan}
Bila data urut, buang separuh daftar di setiap langkah.
\end{Langkah}
\end{column}
\begin{column}{0.31\textwidth}
\begin{Langkah}{3}{Terjemahkan}
Linear search dengan \texttt{for}, binary search dengan \texttt{while} dan tiga indeks.
\end{Langkah}
\end{column}
\end{columns}
\vspace{10pt}
Langkah 1 dan 2 dilatih di Algoritma. Langkah 3 kita kerjakan hari ini dengan C++.\note{Hubungkan eksplisit ke Algoritma minggu ini.}
\end{frame}

\Bagian{2}{Coba dulu, baru dijelaskan}{25 menit eksperimen di GDB Online}
\note{Struggle 25 menit. Berkeliling, jangan menjelaskan materi dulu.}

\begin{frame}{Misi kalian}
\framesubtitle{Struggle, 25 menit}
\begin{columns}[T]
\begin{column}{0.31\textwidth}
\begin{Langkah}{1}{Cari}
Diberikan \texttt{int a[8] = \{4, 9, 2, 7, 5, 9, 1, 8\};}, tampilkan indeks angka 7.
\end{Langkah}
\end{column}
\begin{column}{0.31\textwidth}
\begin{Langkah}{2}{Semua}
Tampilkan semua indeks angka 9, bukan hanya yang pertama.
\end{Langkah}
\end{column}
\begin{column}{0.31\textwidth}
\begin{Langkah}{3}{Tidak ada}
Apa yang ditampilkan program kalian bila mencari angka 6?
\end{Langkah}
\end{column}
\end{columns}
\vspace{8pt}
\begin{columns}[T]
\begin{column}{0.47\textwidth}
\begin{Blok}{Boleh}
Bertanya ke teman, mengubah apa saja, sengaja membuat error.
\end{Blok}
\end{column}
\begin{column}{0.47\textwidth}
\begin{BlokAlert}{Belum dulu}
Mencari jawaban jadi di internet atau AI.
\end{BlokAlert}
\end{column}
\end{columns}
\note{Catat kesalahan yang sering muncul untuk dibahas di debrief.}
\end{frame}

\begin{frame}{Apa yang kalian temukan?}
\framesubtitle{Debrief struggle}
\begin{columns}[T]
\begin{column}{0.31\textwidth}
\begin{BlokGelap}{Pertanyaan 1}
Bagaimana program tahu sebuah data tidak ditemukan?
\end{BlokGelap}
\end{column}
\begin{column}{0.31\textwidth}
\begin{BlokGelap}{Pertanyaan 2}
Kapan \texttt{break} dipakai dan kapan tidak?
\end{BlokGelap}
\end{column}
\begin{column}{0.31\textwidth}
\begin{BlokGelap}{Pertanyaan 3}
Berapa perbandingan untuk mencari elemen terakhir?
\end{BlokGelap}
\end{column}
\end{columns}
\vspace{10pt}
Jawaban kalian akan kita uji di bagian berikutnya.\note{Tulis jawaban di papan; buktikan atau patahkan di bagian materi.}
\end{frame}

\Bagian{3}{Mencari dengan cerdas}{Linear search, binary search, dan menghitung langkah}
\note{Materi 40 menit. Tunjukkan setiap contoh langsung di GDB Online.}

\begin{frame}[fragile]{Linear search}
\framesubtitle{Periksa satu per satu}
\begin{columns}[T]
\begin{column}{0.55\textwidth}
\begin{Kode}[]{linear.cpp}
int a[8] = {4, 9, 2, 7, 5, 9, 1, 8};
int x = 7, posisi = -1;
for (int i = 0; i < 8; i++) {
    if (a[i] == x) {
        posisi = i;
        break;
    }
}
if (posisi == -1) cout << "tidak ada" << endl;
else cout << "indeks " << posisi << endl;
\end{Kode}
\end{column}
\begin{column}{0.41\textwidth}
\only<1>{\begin{TebakDulu}
{\ITERAKickerFont\fontsize{10pt}{12pt}\selectfont\color{ITERAAlert}TEBAK DULU}\par\vspace{4pt}
Sebelum dijalankan, tulis keluarannya di kertas.
\end{TebakDulu}
}\only<2>{\KeluaranFile[]{keluaran/P11/k001.txt}
\begin{Catatan}
Nilai awal -1 menandai belum ditemukan, karena -1 tidak pernah menjadi indeks yang sah.
\end{Catatan}
}
\end{column}
\end{columns}
\note<1>{Minta mahasiswa menebak keluaran sebelum membuka slide berikutnya. Tanyakan satu atau dua tebakan yang berbeda, lalu buktikan.}
\end{frame}

\begin{frame}[fragile]{Mencari semua kemunculan}
\framesubtitle{Tanpa break}
\begin{columns}[T]
\begin{column}{0.60\textwidth}
\begin{Kode}[listing options={style=ITERACodeKecil}]{semua.cpp}
int a[8] = {4, 9, 2, 7, 5, 9, 1, 8};
int x = 9, banyak = 0;
for (int i = 0; i < 8; i++) {
    if (a[i] == x) {
        cout << "indeks " << i << endl;
        banyak++;
    }
}
cout << "muncul " << banyak << " kali" << endl;
\end{Kode}
\end{column}
\begin{column}{0.36\textwidth}
\only<1>{\begin{TebakDulu}
{\ITERAKickerFont\fontsize{10pt}{12pt}\selectfont\color{ITERAAlert}TEBAK DULU}\par\vspace{4pt}
Sebelum dijalankan, tulis keluarannya di kertas.
\end{TebakDulu}
}\only<2>{\KeluaranFile[listing options={style=ITERAPlainKecil}]{keluaran/P11/k002.txt}
\begin{Catatan}
Untuk semua kemunculan, perulangan tidak dihentikan; pencacah mencatat banyaknya.
\end{Catatan}
}
\end{column}
\end{columns}
\note<1>{Minta mahasiswa menebak keluaran sebelum membuka slide berikutnya. Tanyakan satu atau dua tebakan yang berbeda, lalu buktikan.}
\end{frame}

\begin{frame}{Ide binary search}
\framesubtitle{Buang separuh setiap langkah}
Cari 21 pada \{3, 8, 15, 21, 34, 42, 57, 66, 79, 90\}, indeks 0 sampai 9.
\vspace{8pt}

\small
\begin{tabular}{@{}>{\raggedright\arraybackslash}p{110pt}>{\raggedright\arraybackslash}p{82pt}>{\raggedright\arraybackslash}p{82pt}>{\raggedright\arraybackslash}p{82pt}>{\raggedright\arraybackslash}p{110pt}>{\raggedright\arraybackslash}p{358pt}@{}}
\kepala{Langkah} & \kepala{low} & \kepala{high} & \kepala{mid} & \kepala{a[mid]} & \kepala{Keputusan}\\
1 & 0 & 9 & 4 & 34 & 34 > 21, cari di kiri: high = 3\\
\barisgenap 2 & 0 & 3 & 1 & 8 & 8 < 21, cari di kanan: low = 2\\
3 & 2 & 3 & 2 & 15 & 15 < 21, low = 3\\
\barisgenap 4 & 3 & 3 & 3 & 21 & ditemukan di indeks 3\\
\garisdasar
\end{tabular}
\vspace{10pt}

Binary search hanya benar pada data yang sudah terurut.
\end{frame}

\begin{frame}[fragile]{Binary search di C++}
\framesubtitle{low, high, mid}
\begin{columns}[T]
\begin{column}{0.60\textwidth}
\begin{Kode}[listing options={style=ITERACodeKecil}]{biner.cpp}
int a[10] = {3, 8, 15, 21, 34, 42, 57, 66, 79, 90};
int x = 21, low = 0, high = 9, posisi = -1;
while (low <= high) {
    int mid = (low + high) / 2;
    if (a[mid] == x) { posisi = mid; break; }
    else if (a[mid] < x) low = mid + 1;
    else high = mid - 1;
}
cout << "indeks " << posisi << endl;
\end{Kode}
\end{column}
\begin{column}{0.36\textwidth}
\only<1>{\begin{TebakDulu}
{\ITERAKickerFont\fontsize{10pt}{12pt}\selectfont\color{ITERAAlert}TEBAK DULU}\par\vspace{4pt}
Sebelum dijalankan, tulis keluarannya di kertas.
\end{TebakDulu}
}\only<2>{\KeluaranFile[listing options={style=ITERAPlainKecil}]{keluaran/P11/k003.txt}
\begin{Catatan}
Syaratnya \texttt{low <= high}. Dengan \texttt{<}, elemen terakhir yang tersisa tidak sempat diperiksa.
\end{Catatan}
}
\end{column}
\end{columns}
\note<1>{Minta mahasiswa menebak keluaran sebelum membuka slide berikutnya. Tanyakan satu atau dua tebakan yang berbeda, lalu buktikan.}
\end{frame}

\begin{frame}{Berapa langkahnya?}
\framesubtitle{Membandingkan dua algoritma}
\small
\begin{tabular}{@{}>{\raggedright\arraybackslash}p{249pt}>{\raggedright\arraybackslash}p{307pt}>{\raggedright\arraybackslash}p{307pt}@{}}
\kepala{Banyak data} & \kepala{Linear, kasus terburuk} & \kepala{Binary, kasus terburuk}\\
10 & 10 & 4\\
\barisgenap 100 & 100 & 7\\
1.000 & 1.000 & 10\\
\barisgenap 1.000.000 & 1.000.000 & 20\\
\garisdasar
\end{tabular}
\vspace{10pt}

Binary search jauh lebih cepat, tetapi butuh data terurut. Mengurutkan data adalah materi minggu depan.
\end{frame}

\begin{frame}[fragile]{Tebak keluarannya}
\framesubtitle{Latihan menjelaskan, 3 menit}
\begin{columns}[T]
\begin{column}{0.56\textwidth}
\begin{Kode}[listing options={style=ITERACodeKecil}]{tebak.cpp}
int a[7] = {2, 5, 8, 12, 16, 23, 38};
int x = 23, low = 0, high = 6, langkah = 0;
while (low <= high) {
    int mid = (low + high) / 2;
    langkah++;
    if (a[mid] == x) break;
    if (a[mid] < x) low = mid + 1;
    else high = mid - 1;
}
cout << langkah << " " << low << " " << high << endl;
\end{Kode}
\end{column}
\begin{column}{0.40\textwidth}
\begin{Catatan}
Tulis keluarannya di kertas, karakter demi karakter. Jangan dijalankan dulu.
\end{Catatan}
\end{column}
\end{columns}
\note{Contoh soal Bagian B. Beri waktu 3 menit sebelum slide jawaban.}
\end{frame}

\begin{frame}[fragile]{Tebak keluarannya: jawaban}
\framesubtitle{Latihan menjelaskan}
\begin{columns}[T]
\begin{column}{0.56\textwidth}
\begin{Kode}[listing options={style=ITERACodeKecil}]{tebak.cpp}
int a[7] = {2, 5, 8, 12, 16, 23, 38};
int x = 23, low = 0, high = 6, langkah = 0;
while (low <= high) {
    int mid = (low + high) / 2;
    langkah++;
    if (a[mid] == x) break;
    if (a[mid] < x) low = mid + 1;
    else high = mid - 1;
}
cout << langkah << " " << low << " " << high << endl;
\end{Kode}
\end{column}
\begin{column}{0.40\textwidth}
\begin{Keluaran}[]
2 4 6
\end{Keluaran}
\begin{Catatan}
mid = 3 (12 < 23, low = 4), lalu mid = 5 (23, ketemu). Dua langkah; saat berhenti low = 4 dan high = 6.
\end{Catatan}
\end{column}
\end{columns}
\note{Tanyakan jawaban yang paling sering salah, lalu telusuri bersama.}
\end{frame}

\begin{frame}{Error yang sering muncul minggu ini}
\framesubtitle{Pesan diambil langsung dari g++}
\footnotesize
\begin{tabular}{@{}>{\raggedright\arraybackslash}p{208pt}>{\raggedright\arraybackslash}p{256pt}>{\raggedright\arraybackslash}p{398pt}@{}}
\kepala{Kesalahan} & \kepala{Contoh} & \kepala{Pesan pertama dari g++}\\
Membandingkan array dengan == & \texttt{if (a == x)} & \texttt{error: ISO C++ forbids comparison between pointer and integer}\\
\barisgenap Variabel di luar jangkauan & \texttt{mid dipakai di luar while} & \texttt{error: 'mid' was not declared in this scope}\\
Kutip tunggal untuk string & \texttt{nama[i] == 'Rani'} & \texttt{error: no match for 'operator=='}\\
\barisgenap Posisi tidak dipakai & \texttt{int posisi = -1; tak dipakai} & \texttt{warning: unused variable 'posisi'}\\
Kurang titik koma di dalam if & \texttt{posisi = mid break;} & \texttt{error: expected ';' before 'break'}\\
\garisdasar
\end{tabular}
\vspace{10pt}

{\small Pesan di tabel ini dihasilkan dengan g++ dan opsi -Wall, sama seperti di cek.sh.}
\note{Minta mahasiswa memotret slide ini.}
\end{frame}

\Bagian{4}{Hands-on bertingkat}{45 menit, dari dasar sampai tantangan}
\note{Mahasiswa membuka folder 02 Hands-on/P11 - Searching - Hands-on.}

\begin{frame}{Peta hands-on}
\framesubtitle{Folder 02 Hands-on/P11 - Searching - Hands-on}
\small
\begin{tabular}{@{}>{\raggedright\arraybackslash}p{36pt}>{\raggedright\arraybackslash}p{267pt}>{\raggedright\arraybackslash}p{110pt}>{\raggedright\arraybackslash}p{83pt}>{\raggedright\arraybackslash}p{341pt}@{}}
\kepala{No} & \kepala{File} & \kepala{Tingkat} & \kepala{Waktu} & \kepala{Yang dilatih}\\
1 & \texttt{handson1-cari.cpp} & Dasar & 10' & linear search, hasil -1, hitung perbandingan\\
\barisgenap 2 & \texttt{handson2-nama.cpp} & Menengah & 15' & mencari semua kemunculan string\\
3 & \texttt{handson3-biner.cpp} & Tantangan & 20' & trace low-high-mid, perbaiki tiga kesalahan\\
\barisgenap + & \texttt{bonus-banding.cpp} & Pengayaan & sisa & menghitung perbandingan dua algoritma\\
\garisdasar
\end{tabular}
\vspace{10pt}

Kerjakan berurutan. Cocokkan keluaran dengan folder \texttt{expected/}; latihan yang membaca masukan memakai data di folder \texttt{input/}.\note{Saat berkeliling, minta mahasiswa yang mengerjakan latihan tantangan menjelaskan blok PENJELASAN mereka.}
\end{frame}

\begin{frame}{Cara mengecek dan aturannya}
\framesubtitle{Hands-on}
\begin{columns}[T]
\begin{column}{0.47\textwidth}
\begin{Blok}{Di GDB Online}
Salin isi file latihan, lengkapi TODO, klik Run. Bila program membaca masukan, ketik isi file di \texttt{input/}. Bandingkan dengan \texttt{expected/}.
\end{Blok}
\end{column}
\begin{column}{0.47\textwidth}
\begin{BlokGelap}{Di komputer sendiri}
Jalankan \texttt{bash cek.sh}. Skrip mengompilasi dengan \texttt{-Wall -Werror}, memberi masukan dari \texttt{input/}, lalu mencocokkan keluaran.
\end{BlokGelap}
\end{column}
\end{columns}
\vspace{6pt}
\begin{Catatan}
AI boleh untuk memahami pesan error, tidak untuk meminta solusi utuh. Exit Ticket dikerjakan tanpa bantuan apa pun.
\end{Catatan}
\note{Hands-on tidak dinilai langsung; yang dinilai Exit Ticket dan tugas.}
\end{frame}

\Bagian{5}{Exit Ticket dan tugas}{25 menit terakhir, lalu pekerjaan rumah}
\note{Mulai Exit Ticket tepat waktu.}

\begin{frame}{Exit Ticket 11}
\framesubtitle{Individual, 25 menit, tanpa AI dan internet}
\small
\begin{tabular}{@{}>{\raggedright\arraybackslash}p{60pt}>{\raggedright\arraybackslash}p{560pt}>{\raggedright\arraybackslash}p{80pt}>{\raggedright\arraybackslash}p{150pt}@{}}
\kepala{Soal} & \kepala{Isi} & \kepala{Poin} & \kepala{CPMK}\\
A1 & Tulis linear search yang menampilkan semua indeks kemunculan sebuah nilai & 25 & CPMK0607\\
\barisgenap A2 & Tulis binary search pada array terurut yang menampilkan indeks atau tidak ditemukan & 25 & CPMK0607\\
B1 & Isi trace table low, high, mid untuk sebuah pencarian & 20 & CPMK0608\\
\barisgenap B2 & Jelaskan mengapa binary search gagal pada data yang belum terurut & 20 & CPMK0608\\
B3 & Bandingkan banyak perbandingan dua algoritma untuk data yang diberikan & 10 & CPMK0608\\
\garisdasar
\end{tabular}
\vspace{10pt}

{\small Soal A dinilai dari keluaran yang tepat sama. Soal B dinilai dari ketepatan dan alasan. Pelanggaran aturan membuat nilai Exit Ticket ini 0.}
\note{Soal lengkap dibagikan terpisah dan tidak disimpan di repo publik.}
\end{frame}

\begin{frame}{Tugas 11: Direktori Kontak}
\framesubtitle{Dikumpulkan sebelum Pertemuan 12}
\begin{columns}[T]
\begin{column}{0.47\textwidth}
\begin{Blok}{Bagian A, program, 50 poin}
Buat program direktori kontak dengan menu berikut. Data dimasukkan dalam urutan abjad nama; program menolak nama yang tidak urut. Ketentuannya di slide berikut.
\end{Blok}
\end{column}
\begin{column}{0.47\textwidth}
\begin{BlokGelap}{Bagian B, penjelasan, 50 poin}
Komentar maksimal 150 kata: trace table binary search untuk satu nama yang ada dan satu yang tidak ada, serta perbandingan banyak langkah linear dan binary search pada data kalian. Jelaskan mengapa data harus urut abjad.
\end{BlokGelap}
\end{column}
\end{columns}
\vspace{8pt}
{\small Data di tugas adalah data kalian sendiri. Rubrik lengkap ada di Modul Belajar 11.}
\note{Ingatkan bahwa penjelasan disalin dari AI mudah dikenali karena tidak menyebut pengalaman mereka sendiri.}
\end{frame}

\begin{frame}{Ketentuan Tugas 11}
\framesubtitle{Direktori Kontak}
\small
\begin{tabular}{@{}>{\raggedright\arraybackslash}p{87pt}>{\raggedright\arraybackslash}p{788pt}@{}}
\kepala{No} & \kepala{Ketentuan Bagian A}\\
1 & Tambah kontak, maksimal 10 (nama satu kata dan nomor telepon), dengan validasi urutan abjad\\
\barisgenap 2 & Tampilkan semua kontak\\
3 & Cari nama dengan linear search; tampilkan nomor dan banyak perbandingan\\
\barisgenap 4 & Cari nama dengan binary search; tampilkan nomor dan banyak perbandingan\\
5 & Keluar\\
\barisgenap Bonus & Tambahkan pencarian awalan nama, misalnya semua kontak yang diawali "Ra".\\
\garisdasar
\end{tabular}
\note{Bacakan ketentuan satu per satu dan tanyakan apakah ada yang belum jelas.}
\end{frame}

\begin{frame}{Rubrik Tugas 11}
\framesubtitle{Empat level, sama di semua instrumen}
\footnotesize
\begin{tabular}{@{}>{\raggedright\arraybackslash}p{166pt}>{\raggedright\arraybackslash}p{44pt}>{\raggedright\arraybackslash}p{166pt}>{\raggedright\arraybackslash}p{156pt}>{\raggedright\arraybackslash}p{146pt}>{\raggedright\arraybackslash}p{146pt}@{}}
\kepala{Kriteria} & \kepala{Poin} & \kepala{Sangat baik 85--100} & \kepala{Baik 70--84} & \kepala{Cukup 55--69} & \kepala{Kurang <55}\\
A1 Program berjalan & 20 & tanpa error dan peringatan & ada peringatan & perlu satu perbaikan & tidak terkompilasi\\
\barisgenap A2 Kelengkapan menu & 15 & lima menu dan validasi urutan & satu menu kurang & dua kurang & tanpa binary search\\
A3 Ketepatan pencarian & 15 & kedua pencarian tepat, termasuk tidak ditemukan & satu kasus salah & dua salah & pencarian salah\\
\barisgenap B1 Trace table dan perbandingan & 30 & dua trace tepat dan analisis jelas & satu trace keliru & tanpa analisis & tidak ada atau disalin\\
B2 Cerita error dan pengujian & 20 & pesan, sebab, perbaikan, uji & pesan tidak dikutip & hanya perbaikan & tidak ada\\
\garisdasar
\end{tabular}
\vspace{8pt}

{\small Nilai kriteria = poin × skor level. Bagian A total 50, Bagian B total 50.}
\note{Rubrik ini sama strukturnya setiap minggu; yang berubah hanya isi kriteria A2, A3, dan B1.}
\end{frame}

\begin{frame}{Sebelum Pertemuan 12}
\framesubtitle{Persiapan}
\begin{columns}[T]
\begin{column}{0.31\textwidth}
\begin{Langkah}{1}{Selesaikan}
Hands-on yang belum lulus \texttt{cek.sh}, terutama latihan tantangan.
\end{Langkah}
\end{column}
\begin{column}{0.31\textwidth}
\begin{Langkah}{2}{Kumpulkan}
Tugas 11 sebelum Pertemuan 12 dimulai.
\end{Langkah}
\end{column}
\begin{column}{0.31\textwidth}
\begin{Langkah}{3}{Baca}
Modul Belajar 11 untuk mengulang, lalu bagian awal modul berikutnya.
\end{Langkah}
\end{column}
\end{columns}
\vspace{10pt}
Pertemuan 12 membahas sorting: bubble sort dan selection sort, supaya data bisa diurutkan sebelum dicari dengan binary search.\note{Tanyakan satu hal yang paling membingungkan hari ini untuk dibuka di review minggu depan.}
\end{frame}

\SlidePenutup{Terima kasih}{Sampai jumpa di Pertemuan 12: sorting}
\note{Slide penutup.}
\end{document}
"
\documentclass[t]{beamer}
\usetheme{ITERA}
% Mode catatan pembicara: aktif bila \catatan didefinisikan sebelum \input.
\ifdefined\catatan
  \setbeameroption{show only notes}
  \setbeamertemplate{note page}{\vspace*{20pt}\hspace*{4pt}\begin{minipage}{900pt}
    {\ITERAKickerFont\fontsize{11pt}{14pt}\selectfont\color{ITERAGoldText}CATATAN PEMBICARA \ \ |\ \ SLIDE \insertframenumber}\par\vspace{4pt}
    {\ITERADisplay\bfseries\fontsize{22pt}{28pt}\selectfont\color{ITERAStructure}\insertframetitle}\par\vspace{12pt}
    \fontsize{15pt}{23pt}\selectfont\insertnote\end{minipage}}
\fi
\tikzset{fase/.style={draw=none,fill=#1,rounded corners=3pt,minimum width=158pt,minimum height=118pt,text width=146pt,align=center}}

\renewcommand\ITERAmeeting{Pertemuan 11}
\begin{document}
\SlideJudul{IF25-11002 PRAKTIKUM PEMROGRAMAN}{Pertemuan 11\\Searching}{Mencari data di array dengan linear search dan binary search, dan membandingkan banyak langkahnya}{Tim Pengajar}{Tahun 2026. Sejajar dengan Pertemuan 11 Algoritma Pemrograman: Searching.}
\note{Buka dengan menanyakan satu hal dari pertemuan lalu yang masih membingungkan.}

\begin{frame}{Dua mata kuliah, satu minggu}
\framesubtitle{Sebelum mulai}
Topik minggu ini sama di dua mata kuliah, dengan peran yang berbeda.
\vspace{10pt}

\begin{columns}[T]
\begin{column}{0.47\textwidth}
\begin{Blok}{Algoritma: Searching}
Algoritma pencarian beruntun dan biner, prasyarat data terurut, serta menghitung banyak perbandingan.
\end{Blok}
\end{column}
\begin{column}{0.47\textwidth}
\begin{BlokGelap}{Praktikum: Searching}
Menulis linear search dan binary search di C++, menelusuri low, high, dan mid, serta menghitung perbandingan.
\end{BlokGelap}
\end{column}
\end{columns}
\note{Tanyakan apa yang sudah dibahas di Algoritma minggu ini. Gunakan jawabannya sebagai jembatan ke praktikum.}
\end{frame}

\begin{frame}{Saat pulang nanti, kalian bisa}
\framesubtitle{Target hari ini}
\begin{columns}[T]
\begin{column}{0.47\textwidth}
\begin{Langkah}{1}{Menerapkan}
Menulis program yang mencari data pada array dengan linear search dan binary search, termasuk mencari semua kemunculan dan menangani data yang tidak ditemukan

{\footnotesize\color{ITERAInkMuted}Sub-CPMK 11.1, CPMK0607}
\end{Langkah}
\end{column}
\begin{column}{0.47\textwidth}
\begin{Langkah}{2}{Menjelaskan}
Menelusuri langkah binary search dengan trace table low, high, dan mid, serta menjelaskan perbedaan banyak perbandingan kedua algoritma

{\footnotesize\color{ITERAInkMuted}Sub-CPMK 11.2, CPMK0608}
\end{Langkah}
\end{column}
\end{columns}
\vspace{8pt}
\begin{Catatan}
Dua kemampuan ini dinilai sama berat di Exit Ticket, tugas, UTS, dan UAS.
\end{Catatan}
\note{Bacakan target dengan pelan. Kata kuncinya menerapkan dan menjelaskan.}
\end{frame}

\begin{frame}{Ingat minggu lalu}
\framesubtitle{Review, 5 menit}
\begin{columns}[T]
\begin{column}{0.31\textwidth}
\begin{BlokGelap}{Pertanyaan 1}
Bagaimana menghitung rata-rata setiap kolom array 2D?
\end{BlokGelap}
\end{column}
\begin{column}{0.31\textwidth}
\begin{BlokGelap}{Pertanyaan 2}
Apa elemen diagonal utama matriks?
\end{BlokGelap}
\end{column}
\begin{column}{0.31\textwidth}
\begin{BlokGelap}{Pertanyaan 3}
Apa hasil \texttt{s.find("tik")} untuk \texttt{"Praktikum"}?
\end{BlokGelap}
\end{column}
\end{columns}
\vspace{8pt}
Jawab di kertas dulu, lalu cocokkan dengan teman sebelah.\note{Pilih dua mahasiswa untuk menjawab. Koreksi singkat, jangan mengajar ulang.}
\end{frame}

\Bagian{1}{Masalah pembuka}{Mencari nama di daftar wisuda}
\note{Masalah pembuka 10 menit.}

\begin{frame}[fragile]{Mencari nama di daftar wisuda}
\framesubtitle{Masalah minggu ini}
\begin{columns}[T]
\begin{column}{0.55\textwidth}
{Panitia wisuda memegang daftar 1.200 nama yang sudah urut abjad. Seorang orang tua bertanya, apakah anaknya ada di daftar?

\vspace{6pt}
Membaca dari nama pertama sampai ketemu bisa butuh ratusan kali melihat. Padahal saat mencari kata di kamus, kita tidak membaca dari halaman pertama.\par}
\vspace{6pt}
\begin{Catatan}
\textbf{Diskusi 2 menit.} Bagaimana kalian mencari kata di kamus? Mengapa cara itu tidak bisa dipakai pada daftar yang belum urut?
\end{Catatan}
\end{column}
\begin{column}{0.41\textwidth}
\begin{Keluaran}[]
Cari: Siti Rahma
Ditemukan di nomor 1023
Linear search : 1023 perbandingan
Binary search : 10 perbandingan
\end{Keluaran}
\end{column}
\end{columns}
\note{Tampung beberapa jawaban tanpa menilai benar salah.}
\end{frame}

\begin{frame}{Dari langkah ke instruksi}
\framesubtitle{Sambungan dengan Algoritma}
\begin{columns}[T]
\begin{column}{0.31\textwidth}
\begin{Langkah}{1}{Bandingkan satu per satu}
Cara paling sederhana: periksa setiap elemen dari depan.
\end{Langkah}
\end{column}
\begin{column}{0.31\textwidth}
\begin{Langkah}{2}{Manfaatkan urutan}
Bila data urut, buang separuh daftar di setiap langkah.
\end{Langkah}
\end{column}
\begin{column}{0.31\textwidth}
\begin{Langkah}{3}{Terjemahkan}
Linear search dengan \texttt{for}, binary search dengan \texttt{while} dan tiga indeks.
\end{Langkah}
\end{column}
\end{columns}
\vspace{10pt}
Langkah 1 dan 2 dilatih di Algoritma. Langkah 3 kita kerjakan hari ini dengan C++.\note{Hubungkan eksplisit ke Algoritma minggu ini.}
\end{frame}

\Bagian{2}{Coba dulu, baru dijelaskan}{25 menit eksperimen di GDB Online}
\note{Struggle 25 menit. Berkeliling, jangan menjelaskan materi dulu.}

\begin{frame}{Misi kalian}
\framesubtitle{Struggle, 25 menit}
\begin{columns}[T]
\begin{column}{0.31\textwidth}
\begin{Langkah}{1}{Cari}
Diberikan \texttt{int a[8] = \{4, 9, 2, 7, 5, 9, 1, 8\};}, tampilkan indeks angka 7.
\end{Langkah}
\end{column}
\begin{column}{0.31\textwidth}
\begin{Langkah}{2}{Semua}
Tampilkan semua indeks angka 9, bukan hanya yang pertama.
\end{Langkah}
\end{column}
\begin{column}{0.31\textwidth}
\begin{Langkah}{3}{Tidak ada}
Apa yang ditampilkan program kalian bila mencari angka 6?
\end{Langkah}
\end{column}
\end{columns}
\vspace{8pt}
\begin{columns}[T]
\begin{column}{0.47\textwidth}
\begin{Blok}{Boleh}
Bertanya ke teman, mengubah apa saja, sengaja membuat error.
\end{Blok}
\end{column}
\begin{column}{0.47\textwidth}
\begin{BlokAlert}{Belum dulu}
Mencari jawaban jadi di internet atau AI.
\end{BlokAlert}
\end{column}
\end{columns}
\note{Catat kesalahan yang sering muncul untuk dibahas di debrief.}
\end{frame}

\begin{frame}{Apa yang kalian temukan?}
\framesubtitle{Debrief struggle}
\begin{columns}[T]
\begin{column}{0.31\textwidth}
\begin{BlokGelap}{Pertanyaan 1}
Bagaimana program tahu sebuah data tidak ditemukan?
\end{BlokGelap}
\end{column}
\begin{column}{0.31\textwidth}
\begin{BlokGelap}{Pertanyaan 2}
Kapan \texttt{break} dipakai dan kapan tidak?
\end{BlokGelap}
\end{column}
\begin{column}{0.31\textwidth}
\begin{BlokGelap}{Pertanyaan 3}
Berapa perbandingan untuk mencari elemen terakhir?
\end{BlokGelap}
\end{column}
\end{columns}
\vspace{10pt}
Jawaban kalian akan kita uji di bagian berikutnya.\note{Tulis jawaban di papan; buktikan atau patahkan di bagian materi.}
\end{frame}

\Bagian{3}{Mencari dengan cerdas}{Linear search, binary search, dan menghitung langkah}
\note{Materi 40 menit. Tunjukkan setiap contoh langsung di GDB Online.}

\begin{frame}[fragile]{Linear search}
\framesubtitle{Periksa satu per satu}
\begin{columns}[T]
\begin{column}{0.55\textwidth}
\begin{Kode}[]{linear.cpp}
int a[8] = {4, 9, 2, 7, 5, 9, 1, 8};
int x = 7, posisi = -1;
for (int i = 0; i < 8; i++) {
    if (a[i] == x) {
        posisi = i;
        break;
    }
}
if (posisi == -1) cout << "tidak ada" << endl;
else cout << "indeks " << posisi << endl;
\end{Kode}
\end{column}
\begin{column}{0.41\textwidth}
\only<1>{\begin{TebakDulu}
{\ITERAKickerFont\fontsize{10pt}{12pt}\selectfont\color{ITERAAlert}TEBAK DULU}\par\vspace{4pt}
Sebelum dijalankan, tulis keluarannya di kertas.
\end{TebakDulu}
}\only<2>{\KeluaranFile[]{keluaran/P11/k001.txt}
\begin{Catatan}
Nilai awal -1 menandai belum ditemukan, karena -1 tidak pernah menjadi indeks yang sah.
\end{Catatan}
}
\end{column}
\end{columns}
\note<1>{Minta mahasiswa menebak keluaran sebelum membuka slide berikutnya. Tanyakan satu atau dua tebakan yang berbeda, lalu buktikan.}
\end{frame}

\begin{frame}[fragile]{Mencari semua kemunculan}
\framesubtitle{Tanpa break}
\begin{columns}[T]
\begin{column}{0.60\textwidth}
\begin{Kode}[listing options={style=ITERACodeKecil}]{semua.cpp}
int a[8] = {4, 9, 2, 7, 5, 9, 1, 8};
int x = 9, banyak = 0;
for (int i = 0; i < 8; i++) {
    if (a[i] == x) {
        cout << "indeks " << i << endl;
        banyak++;
    }
}
cout << "muncul " << banyak << " kali" << endl;
\end{Kode}
\end{column}
\begin{column}{0.36\textwidth}
\only<1>{\begin{TebakDulu}
{\ITERAKickerFont\fontsize{10pt}{12pt}\selectfont\color{ITERAAlert}TEBAK DULU}\par\vspace{4pt}
Sebelum dijalankan, tulis keluarannya di kertas.
\end{TebakDulu}
}\only<2>{\KeluaranFile[listing options={style=ITERAPlainKecil}]{keluaran/P11/k002.txt}
\begin{Catatan}
Untuk semua kemunculan, perulangan tidak dihentikan; pencacah mencatat banyaknya.
\end{Catatan}
}
\end{column}
\end{columns}
\note<1>{Minta mahasiswa menebak keluaran sebelum membuka slide berikutnya. Tanyakan satu atau dua tebakan yang berbeda, lalu buktikan.}
\end{frame}

\begin{frame}{Ide binary search}
\framesubtitle{Buang separuh setiap langkah}
Cari 21 pada \{3, 8, 15, 21, 34, 42, 57, 66, 79, 90\}, indeks 0 sampai 9.
\vspace{8pt}

\small
\begin{tabular}{@{}>{\raggedright\arraybackslash}p{110pt}>{\raggedright\arraybackslash}p{82pt}>{\raggedright\arraybackslash}p{82pt}>{\raggedright\arraybackslash}p{82pt}>{\raggedright\arraybackslash}p{110pt}>{\raggedright\arraybackslash}p{358pt}@{}}
\kepala{Langkah} & \kepala{low} & \kepala{high} & \kepala{mid} & \kepala{a[mid]} & \kepala{Keputusan}\\
1 & 0 & 9 & 4 & 34 & 34 > 21, cari di kiri: high = 3\\
\barisgenap 2 & 0 & 3 & 1 & 8 & 8 < 21, cari di kanan: low = 2\\
3 & 2 & 3 & 2 & 15 & 15 < 21, low = 3\\
\barisgenap 4 & 3 & 3 & 3 & 21 & ditemukan di indeks 3\\
\garisdasar
\end{tabular}
\vspace{10pt}

Binary search hanya benar pada data yang sudah terurut.
\end{frame}

\begin{frame}[fragile]{Binary search di C++}
\framesubtitle{low, high, mid}
\begin{columns}[T]
\begin{column}{0.60\textwidth}
\begin{Kode}[listing options={style=ITERACodeKecil}]{biner.cpp}
int a[10] = {3, 8, 15, 21, 34, 42, 57, 66, 79, 90};
int x = 21, low = 0, high = 9, posisi = -1;
while (low <= high) {
    int mid = (low + high) / 2;
    if (a[mid] == x) { posisi = mid; break; }
    else if (a[mid] < x) low = mid + 1;
    else high = mid - 1;
}
cout << "indeks " << posisi << endl;
\end{Kode}
\end{column}
\begin{column}{0.36\textwidth}
\only<1>{\begin{TebakDulu}
{\ITERAKickerFont\fontsize{10pt}{12pt}\selectfont\color{ITERAAlert}TEBAK DULU}\par\vspace{4pt}
Sebelum dijalankan, tulis keluarannya di kertas.
\end{TebakDulu}
}\only<2>{\KeluaranFile[listing options={style=ITERAPlainKecil}]{keluaran/P11/k003.txt}
\begin{Catatan}
Syaratnya \texttt{low <= high}. Dengan \texttt{<}, elemen terakhir yang tersisa tidak sempat diperiksa.
\end{Catatan}
}
\end{column}
\end{columns}
\note<1>{Minta mahasiswa menebak keluaran sebelum membuka slide berikutnya. Tanyakan satu atau dua tebakan yang berbeda, lalu buktikan.}
\end{frame}

\begin{frame}{Berapa langkahnya?}
\framesubtitle{Membandingkan dua algoritma}
\small
\begin{tabular}{@{}>{\raggedright\arraybackslash}p{249pt}>{\raggedright\arraybackslash}p{307pt}>{\raggedright\arraybackslash}p{307pt}@{}}
\kepala{Banyak data} & \kepala{Linear, kasus terburuk} & \kepala{Binary, kasus terburuk}\\
10 & 10 & 4\\
\barisgenap 100 & 100 & 7\\
1.000 & 1.000 & 10\\
\barisgenap 1.000.000 & 1.000.000 & 20\\
\garisdasar
\end{tabular}
\vspace{10pt}

Binary search jauh lebih cepat, tetapi butuh data terurut. Mengurutkan data adalah materi minggu depan.
\end{frame}

\begin{frame}[fragile]{Tebak keluarannya}
\framesubtitle{Latihan menjelaskan, 3 menit}
\begin{columns}[T]
\begin{column}{0.56\textwidth}
\begin{Kode}[listing options={style=ITERACodeKecil}]{tebak.cpp}
int a[7] = {2, 5, 8, 12, 16, 23, 38};
int x = 23, low = 0, high = 6, langkah = 0;
while (low <= high) {
    int mid = (low + high) / 2;
    langkah++;
    if (a[mid] == x) break;
    if (a[mid] < x) low = mid + 1;
    else high = mid - 1;
}
cout << langkah << " " << low << " " << high << endl;
\end{Kode}
\end{column}
\begin{column}{0.40\textwidth}
\begin{Catatan}
Tulis keluarannya di kertas, karakter demi karakter. Jangan dijalankan dulu.
\end{Catatan}
\end{column}
\end{columns}
\note{Contoh soal Bagian B. Beri waktu 3 menit sebelum slide jawaban.}
\end{frame}

\begin{frame}[fragile]{Tebak keluarannya: jawaban}
\framesubtitle{Latihan menjelaskan}
\begin{columns}[T]
\begin{column}{0.56\textwidth}
\begin{Kode}[listing options={style=ITERACodeKecil}]{tebak.cpp}
int a[7] = {2, 5, 8, 12, 16, 23, 38};
int x = 23, low = 0, high = 6, langkah = 0;
while (low <= high) {
    int mid = (low + high) / 2;
    langkah++;
    if (a[mid] == x) break;
    if (a[mid] < x) low = mid + 1;
    else high = mid - 1;
}
cout << langkah << " " << low << " " << high << endl;
\end{Kode}
\end{column}
\begin{column}{0.40\textwidth}
\begin{Keluaran}[]
2 4 6
\end{Keluaran}
\begin{Catatan}
mid = 3 (12 < 23, low = 4), lalu mid = 5 (23, ketemu). Dua langkah; saat berhenti low = 4 dan high = 6.
\end{Catatan}
\end{column}
\end{columns}
\note{Tanyakan jawaban yang paling sering salah, lalu telusuri bersama.}
\end{frame}

\begin{frame}{Error yang sering muncul minggu ini}
\framesubtitle{Pesan diambil langsung dari g++}
\footnotesize
\begin{tabular}{@{}>{\raggedright\arraybackslash}p{208pt}>{\raggedright\arraybackslash}p{256pt}>{\raggedright\arraybackslash}p{398pt}@{}}
\kepala{Kesalahan} & \kepala{Contoh} & \kepala{Pesan pertama dari g++}\\
Membandingkan array dengan == & \texttt{if (a == x)} & \texttt{error: ISO C++ forbids comparison between pointer and integer}\\
\barisgenap Variabel di luar jangkauan & \texttt{mid dipakai di luar while} & \texttt{error: 'mid' was not declared in this scope}\\
Kutip tunggal untuk string & \texttt{nama[i] == 'Rani'} & \texttt{error: no match for 'operator=='}\\
\barisgenap Posisi tidak dipakai & \texttt{int posisi = -1; tak dipakai} & \texttt{warning: unused variable 'posisi'}\\
Kurang titik koma di dalam if & \texttt{posisi = mid break;} & \texttt{error: expected ';' before 'break'}\\
\garisdasar
\end{tabular}
\vspace{10pt}

{\small Pesan di tabel ini dihasilkan dengan g++ dan opsi -Wall, sama seperti di cek.sh.}
\note{Minta mahasiswa memotret slide ini.}
\end{frame}

\Bagian{4}{Hands-on bertingkat}{45 menit, dari dasar sampai tantangan}
\note{Mahasiswa membuka folder 02 Hands-on/P11 - Searching - Hands-on.}

\begin{frame}{Peta hands-on}
\framesubtitle{Folder 02 Hands-on/P11 - Searching - Hands-on}
\small
\begin{tabular}{@{}>{\raggedright\arraybackslash}p{36pt}>{\raggedright\arraybackslash}p{267pt}>{\raggedright\arraybackslash}p{110pt}>{\raggedright\arraybackslash}p{83pt}>{\raggedright\arraybackslash}p{341pt}@{}}
\kepala{No} & \kepala{File} & \kepala{Tingkat} & \kepala{Waktu} & \kepala{Yang dilatih}\\
1 & \texttt{handson1-cari.cpp} & Dasar & 10' & linear search, hasil -1, hitung perbandingan\\
\barisgenap 2 & \texttt{handson2-nama.cpp} & Menengah & 15' & mencari semua kemunculan string\\
3 & \texttt{handson3-biner.cpp} & Tantangan & 20' & trace low-high-mid, perbaiki tiga kesalahan\\
\barisgenap + & \texttt{bonus-banding.cpp} & Pengayaan & sisa & menghitung perbandingan dua algoritma\\
\garisdasar
\end{tabular}
\vspace{10pt}

Kerjakan berurutan. Cocokkan keluaran dengan folder \texttt{expected/}; latihan yang membaca masukan memakai data di folder \texttt{input/}.\note{Saat berkeliling, minta mahasiswa yang mengerjakan latihan tantangan menjelaskan blok PENJELASAN mereka.}
\end{frame}

\begin{frame}{Cara mengecek dan aturannya}
\framesubtitle{Hands-on}
\begin{columns}[T]
\begin{column}{0.47\textwidth}
\begin{Blok}{Di GDB Online}
Salin isi file latihan, lengkapi TODO, klik Run. Bila program membaca masukan, ketik isi file di \texttt{input/}. Bandingkan dengan \texttt{expected/}.
\end{Blok}
\end{column}
\begin{column}{0.47\textwidth}
\begin{BlokGelap}{Di komputer sendiri}
Jalankan \texttt{bash cek.sh}. Skrip mengompilasi dengan \texttt{-Wall -Werror}, memberi masukan dari \texttt{input/}, lalu mencocokkan keluaran.
\end{BlokGelap}
\end{column}
\end{columns}
\vspace{6pt}
\begin{Catatan}
AI boleh untuk memahami pesan error, tidak untuk meminta solusi utuh. Exit Ticket dikerjakan tanpa bantuan apa pun.
\end{Catatan}
\note{Hands-on tidak dinilai langsung; yang dinilai Exit Ticket dan tugas.}
\end{frame}

\Bagian{5}{Exit Ticket dan tugas}{25 menit terakhir, lalu pekerjaan rumah}
\note{Mulai Exit Ticket tepat waktu.}

\begin{frame}{Exit Ticket 11}
\framesubtitle{Individual, 25 menit, tanpa AI dan internet}
\small
\begin{tabular}{@{}>{\raggedright\arraybackslash}p{60pt}>{\raggedright\arraybackslash}p{560pt}>{\raggedright\arraybackslash}p{80pt}>{\raggedright\arraybackslash}p{150pt}@{}}
\kepala{Soal} & \kepala{Isi} & \kepala{Poin} & \kepala{CPMK}\\
A1 & Tulis linear search yang menampilkan semua indeks kemunculan sebuah nilai & 25 & CPMK0607\\
\barisgenap A2 & Tulis binary search pada array terurut yang menampilkan indeks atau tidak ditemukan & 25 & CPMK0607\\
B1 & Isi trace table low, high, mid untuk sebuah pencarian & 20 & CPMK0608\\
\barisgenap B2 & Jelaskan mengapa binary search gagal pada data yang belum terurut & 20 & CPMK0608\\
B3 & Bandingkan banyak perbandingan dua algoritma untuk data yang diberikan & 10 & CPMK0608\\
\garisdasar
\end{tabular}
\vspace{10pt}

{\small Soal A dinilai dari keluaran yang tepat sama. Soal B dinilai dari ketepatan dan alasan. Pelanggaran aturan membuat nilai Exit Ticket ini 0.}
\note{Soal lengkap dibagikan terpisah dan tidak disimpan di repo publik.}
\end{frame}

\begin{frame}{Tugas 11: Direktori Kontak}
\framesubtitle{Dikumpulkan sebelum Pertemuan 12}
\begin{columns}[T]
\begin{column}{0.47\textwidth}
\begin{Blok}{Bagian A, program, 50 poin}
Buat program direktori kontak dengan menu berikut. Data dimasukkan dalam urutan abjad nama; program menolak nama yang tidak urut. Ketentuannya di slide berikut.
\end{Blok}
\end{column}
\begin{column}{0.47\textwidth}
\begin{BlokGelap}{Bagian B, penjelasan, 50 poin}
Komentar maksimal 150 kata: trace table binary search untuk satu nama yang ada dan satu yang tidak ada, serta perbandingan banyak langkah linear dan binary search pada data kalian. Jelaskan mengapa data harus urut abjad.
\end{BlokGelap}
\end{column}
\end{columns}
\vspace{8pt}
{\small Data di tugas adalah data kalian sendiri. Rubrik lengkap ada di Modul Belajar 11.}
\note{Ingatkan bahwa penjelasan disalin dari AI mudah dikenali karena tidak menyebut pengalaman mereka sendiri.}
\end{frame}

\begin{frame}{Ketentuan Tugas 11}
\framesubtitle{Direktori Kontak}
\small
\begin{tabular}{@{}>{\raggedright\arraybackslash}p{87pt}>{\raggedright\arraybackslash}p{788pt}@{}}
\kepala{No} & \kepala{Ketentuan Bagian A}\\
1 & Tambah kontak, maksimal 10 (nama satu kata dan nomor telepon), dengan validasi urutan abjad\\
\barisgenap 2 & Tampilkan semua kontak\\
3 & Cari nama dengan linear search; tampilkan nomor dan banyak perbandingan\\
\barisgenap 4 & Cari nama dengan binary search; tampilkan nomor dan banyak perbandingan\\
5 & Keluar\\
\barisgenap Bonus & Tambahkan pencarian awalan nama, misalnya semua kontak yang diawali "Ra".\\
\garisdasar
\end{tabular}
\note{Bacakan ketentuan satu per satu dan tanyakan apakah ada yang belum jelas.}
\end{frame}

\begin{frame}{Rubrik Tugas 11}
\framesubtitle{Empat level, sama di semua instrumen}
\footnotesize
\begin{tabular}{@{}>{\raggedright\arraybackslash}p{166pt}>{\raggedright\arraybackslash}p{44pt}>{\raggedright\arraybackslash}p{166pt}>{\raggedright\arraybackslash}p{156pt}>{\raggedright\arraybackslash}p{146pt}>{\raggedright\arraybackslash}p{146pt}@{}}
\kepala{Kriteria} & \kepala{Poin} & \kepala{Sangat baik 85--100} & \kepala{Baik 70--84} & \kepala{Cukup 55--69} & \kepala{Kurang <55}\\
A1 Program berjalan & 20 & tanpa error dan peringatan & ada peringatan & perlu satu perbaikan & tidak terkompilasi\\
\barisgenap A2 Kelengkapan menu & 15 & lima menu dan validasi urutan & satu menu kurang & dua kurang & tanpa binary search\\
A3 Ketepatan pencarian & 15 & kedua pencarian tepat, termasuk tidak ditemukan & satu kasus salah & dua salah & pencarian salah\\
\barisgenap B1 Trace table dan perbandingan & 30 & dua trace tepat dan analisis jelas & satu trace keliru & tanpa analisis & tidak ada atau disalin\\
B2 Cerita error dan pengujian & 20 & pesan, sebab, perbaikan, uji & pesan tidak dikutip & hanya perbaikan & tidak ada\\
\garisdasar
\end{tabular}
\vspace{8pt}

{\small Nilai kriteria = poin × skor level. Bagian A total 50, Bagian B total 50.}
\note{Rubrik ini sama strukturnya setiap minggu; yang berubah hanya isi kriteria A2, A3, dan B1.}
\end{frame}

\begin{frame}{Sebelum Pertemuan 12}
\framesubtitle{Persiapan}
\begin{columns}[T]
\begin{column}{0.31\textwidth}
\begin{Langkah}{1}{Selesaikan}
Hands-on yang belum lulus \texttt{cek.sh}, terutama latihan tantangan.
\end{Langkah}
\end{column}
\begin{column}{0.31\textwidth}
\begin{Langkah}{2}{Kumpulkan}
Tugas 11 sebelum Pertemuan 12 dimulai.
\end{Langkah}
\end{column}
\begin{column}{0.31\textwidth}
\begin{Langkah}{3}{Baca}
Modul Belajar 11 untuk mengulang, lalu bagian awal modul berikutnya.
\end{Langkah}
\end{column}
\end{columns}
\vspace{10pt}
Pertemuan 12 membahas sorting: bubble sort dan selection sort, supaya data bisa diurutkan sebelum dicari dengan binary search.\note{Tanyakan satu hal yang paling membingungkan hari ini untuk dibuka di review minggu depan.}
\end{frame}

\SlidePenutup{Terima kasih}{Sampai jumpa di Pertemuan 12: sorting}
\note{Slide penutup.}
\end{document}
